\documentclass[10pt,a4paper]{article}

% ----------------- Paquetes -------------------%
\usepackage[T1]{fontenc}
\usepackage[utf8]{inputenc}
\usepackage[a4paper,margin=2.5cm]{geometry}
\usepackage{mathptmx}             
\usepackage{changepage} 
\usepackage{setspace}
\usepackage[spanish]{babel}
\usepackage[labelfont=bf,labelsep=space]{caption}
\usepackage{amssymb}
\usepackage{amsmath}
\usepackage{ebgaramond}
\usepackage{placeins}
\usepackage{etoolbox}
\usepackage{setspace}
\usepackage{parskip} 
\usepackage{tcolorbox}
\usepackage{needspace}
\usepackage{xparse}
\usepackage{ocgx2}
\usepackage{hyperref}
\usepackage{hypcap}



%%%%%%%%%%%%%%%%%%%%%%%%%%%%%%%%%%%%%%%%%%%%%%%%%%%%%%%%%%%%%%%%
% --- Fija primero tus valores globales "buenos" ---
\setstretch{1.2}                 % Interlineado global
\setlength{\parskip}{0.7\baselineskip}
\setlength{\parindent}{0pt}
\newlength{\GlobalParskip}
\setlength{\GlobalParskip}{\parskip}

\newcounter{eqnum}[subsection]
\renewcommand{\theeqnum}{\arabic{eqnum}}
\newcounter{alphaitem}
\newcounter{numitem}

%% ------------------- Recuadros----------------------%%
\newcommand{\puntosclave}[1]{%
  \begin{tcolorbox}[
    colframe=black,
    title={Puntos clave},
    before upper={\setlength{\parindent}{0.5cm}\setlength{\parskip}{0.5\baselineskip}}
  ]
  #1
  \end{tcolorbox}
}

\newcommand{\modificaciones}[1]{
  \begin{tcolorbox}[
    colback=white,
    colframe=red,
    title={Modificaciones a realizar},
    before upper={\setlength{\parindent}{0.5cm}\setlength{\parskip}{0.5\baselineskip}}
  ]
  #1
  \end{tcolorbox}
}

%% ------------------- Preguntas TEST----------------------%%
% Opciones: radio (single-choice)
\NewDocumentCommand{\OptRadio}{m m m}{%
  \noindent\CheckBox[radio,name=#1,value=#2,width=1.2ex,height=1.2ex]{}%
  \;\textbf{#2.}~#3\par
}

% Opciones: checkbox (multi-choice)
\NewDocumentCommand{\OptCheck}{m m}{%
  \noindent\CheckBox[name=#1,width=1.2ex,height=1.2ex,bordercolor={0 0 0}]{}%
  \;#2\par
}

% Pregunta single-choice (radio)
% Args: 1=id, 2=enunciado, 3=opciones, 4=solución+justificación, 5=imagen (o {}), 6=origen
\NewDocumentCommand{\PreguntaTestSingle}{m m m m m m}{%
  \par\medskip
  \noindent\textbf{#1.}~#2\par\smallskip

    % Imagen (si no es {})
    \IfBlankTF{#5}{}{%
      \begin{center}
        \includegraphics[width=0.75\linewidth]{#5}
      \end{center}
    }

    \begin{Form}
      #3
    \end{Form}

    \par\smallskip
    \switchocg{sol-#1}{\fbox{Mostrar / ocultar solución}}%
    \hfill{\footnotesize #6}\par\smallskip

    \begin{ocg}{Solución}{sol-#1}{0}
      \noindent #4
    \end{ocg}
  \par\medskip
}

% Pregunta multi-choice (checkboxes)
\NewDocumentCommand{\PreguntaTestMulti}{m m m m m m}{%
  \par\medskip
  \noindent\textbf{#1.}~#2\par\smallskip

    % Imagen (si no es {})
    \IfBlankTF{#5}{}{%
      \begin{center}
        \includegraphics[width=0.75\linewidth]{#5}
      \end{center}
    }
    \begin{Form}
      #3
    \end{Form}

    \par\smallskip
    \switchocg{sol-#1}{\fbox{Mostrar / ocultar solución}}%
    \hfill{\footnotesize #6}\par\smallskip

    \begin{ocg}{Solución}{sol-#1}{0}
      \noindent #4
    \end{ocg}
  \par\medskip
}

%% ------------------- Listas----------------------%%
    % --- Lista alfabética ---
    \newenvironment{la}
      {\begin{list}{\alph{alphaitem})}{%
          \usecounter{alphaitem}%
          \setlength{\leftmargin}{1cm}%
          \setlength{\parskip}{3pt}% 
          \setlength{\itemsep}{0pt}% ← espacio entre ítems
          \setlength{\parsep}{0pt}%  ← párrafos dentro de un ítem
      }}
      {\end{list}}
      
    % --- Lista numérica ---
    \newenvironment{lnum}
      {\begin{list}{\arabic{numitem}.}{%
          \usecounter{numitem}%
          \setlength{\leftmargin}{1cm}%
          \setlength{\parskip}{3pt}% 
          \setlength{\itemsep}{0pt}% ← espacio entre ítems
          \setlength{\parsep}{0pt}%  ← párrafos dentro de un ítem
      }}
      {\end{list}}

%% ------------------- Preguntas Test ----------------------%%
      \usepackage{xparse} % si ya lo tienes, no lo repitas

    \NewDocumentEnvironment{test}{m m o}{%
      \begin{tcolorbox}[
        colback=white,
        colframe=black!15,
        boxrule=0.6pt,
        arc=2mm,
        left=2mm,right=2mm,top=2mm,bottom=2mm
      ]
      \centering
      \includegraphics[width=\linewidth]{#1}
      \end{tcolorbox}
    
      \vspace{2mm}
    
      % Caja SOLO para texto (SÍ partible y mantiene márgenes al partir)
      \begin{tcolorbox}[
        enhanced,
        breakable,
        colback=black!3,
        colframe=black!25,
        boxrule=0.6pt,
        arc=2mm,
        left=2mm,right=2mm,top=1.5mm,bottom=1.5mm
      ]
      \textbf{Respuesta:} \textbf{#2}%
      \IfValueT{#3}{\par\smallskip \textbf{Justificación:} #3}
      \end{tcolorbox}
    }{}
      
% ------------------- Imágenes----------------------%
    \graphicspath{{Img/}}% Rutas por defecto 
    % --- Imagen adaptativa ---
    % Registros de dimensión definidos una sola vez
    \newdimen\imgcap
    \newdimen\imgmin
    \newdimen\imgmax
    \newdimen\imgremain
    \newdimen\imgh
    \newdimen\imgneed
    
    \newcommand{\imagenfit}[3]{%
      \addvspace{10pt}% separación antes
      \par\begingroup
        % Parámetros de composición (asignaciones, no \newdimen)
        \imgcap=1\baselineskip            % alto estimado de título+fuente
        \imgmin=0.15\textheight
        \imgmax=0.15\textheight
        \imgremain=\pagegoal
        \ifdim\pagegoal=\maxdimen \imgremain=0pt\fi
        \advance\imgremain by -\pagetotal
        \advance\imgremain by -\imgcap
        % Decisión de altura destino
        \ifdim\imgremain<\imgmin
          \newpage
          \imgh=0.20\textheight
        \else
          \imgh=\imgremain
          \ifdim\imgh>\imgmax \imgh=\imgmax \fi
        \fi
        % Reservar espacio para evitar cortes del bloque completo
        \imgneed=\imgh \advance\imgneed by \imgcap
        \Needspace{\imgneed}
        % Cuerpo
        \noindent\begin{minipage}{\textwidth}
          \centering
          \captionof{figure}{\textemdash\ #2}\par\vspace{0.3em}% título arriba
          \includegraphics[width=\linewidth,height=\imgh,keepaspectratio]{\detokenize{#1}}\par
          \vspace{0.3em}\small\textit{Fuente: }#3% fuente abajo
        \end{minipage}%
      \endgroup
      \par\addvspace{10pt}%
    }

%------------ Bloque de RECUADRO ESPECIAL -------------%
    \newenvironment{specialblock}[1]{%
      \begin{quote} % crea sangría a izquierda y derecha
      \small % texto más pequeño
      \noindent\textbf{#1}\\[-0.3em] % título en negrita
      \rule{\linewidth}{0.4pt}\\[0.6em]}{
      \end{quote}}
      
%-------------------------- Portada -----------------------%
\newcommand{\oldtitle}[1]{\def\OldTitle{#1}}
\newcommand{\changerationale}[1]{\def\ChangeWhy{#1}}

\newcommand{\changebox}{%
\begin{tcolorbox}[title={Historial de cambios del tema}]
\textbf{Título anterior:} \OldTitle\par
\textbf{Motivo del cambio:} \ChangeWhy
\end{tcolorbox}
}

%-------------------------- Author Font -----------------------%
    \newcommand{\authorfont}[1]{{\fontfamily{EBGaramond-TLF}\selectfont #1}}

%-------------------------- Anexos -----------------------%
    \makeatletter
    \newcounter{anexo}
    \renewcommand{\theanexo}{\Roman{anexo}}
    \newcommand{\anexo}[1]{%
      \clearpage
      \phantomsection
      \refstepcounter{anexo}%
      \noindent\textbf{Anexo \theanexo.- }#1\par\medskip
      \addcontentsline{stc}{section}{Anexo \theanexo.- #1}%
    }
    \makeatother

    %-------------------------- Lista de figuras (personal) -----------------------%
\makeatletter
\newcommand{\MyListOfFigures}{%
  \clearpage
  \phantomsection
  {\centering\bfseries\Large Índice de figuras\par}%
  \medskip
  % Entrada en nuestro índice de contenido (stc)
  \addcontentsline{stc}{section}{Índice de figuras}%
  % Recuperación del listado (sin cabecera propia)
  \begingroup
    \parindent=0pt\parskip=2pt
    \@starttoc{lof}%
  \endgroup
}
\makeatother

%-------------------------- Bibliografía -----------------------%
\makeatletter
% Contador para las referencias
\newcounter{myref}
% Cabecera de la bibliografía, entra en el índice 'stc'
\newcommand{\MyBibliography}{%
  \clearpage
  \phantomsection
  {\centering\bfseries\Large Bibliografía y referencias\par}%
  \medskip
  \addcontentsline{stc}{section}{Bibliografía y referencias}%
  \setcounter{myref}{0}%
}
\newcommand{\MyBibItem}[4]{%
  \refstepcounter{myref}%
  \noindent[\themyref]\ #1.\ \emph{#2}. #3. #4\par
}
\makeatother

%--------- Establecer cuatro niveles de índice --------%
    \setcounter{tocdepth}{4}   
    \setcounter{secnumdepth}{4}% numera hasta \paragraph

%%%%%%%%%%%%%%%%%%%%%%%%%%%%%%%%

\makeatletter

    %--------- Interlineado Global --------%
    \edef\GlobalBaseLineStretch{\baselinestretch}
    
    %--------- Espaciado de seccinones --------%
    \renewcommand\section{\@startsection{section}
      {1}{\z@}
      {2.5ex \@plus 1ex \@minus .2ex} % espacio antes
      {1.5ex \@plus .2ex}             % espacio después
      {\normalfont\Large\bfseries}    % formato del título
    }
    \renewcommand\subsection{\@startsection{subsection}{2}{\z@}
      {3ex \@plus 0.8ex \@minus .2ex}
      {1.5ex \@plus .2ex}
      {\normalfont\large\bfseries}
    }
    \renewcommand\subsubsection{\@startsection{subsubsection}{3}{\z@}
      {3ex \@plus 0.5ex \@minus .2ex}
      {1ex \@plus .2ex}
      {\normalfont\normalsize\bfseries}
    }
    \renewcommand\paragraph{\@startsection{paragraph}{4}{\z@}%
    {-2ex \@plus -0.5ex \@minus -.2ex}% espacio antes
    {0.8ex \@plus .2ex}% espacio después
    {\normalfont\normalsize\itshape}}% ← cursiva en lugar de negrita

    %--------- Ecuaciones --------%   
    % Variables globales para almacenar títulos y notas
    \gdef\leq@current@section{}%
    \gdef\leq@current@subsection{}%
    \gdef\leq@last@printed@section{}%
    \gdef\leq@last@printed@subsection{}%
    
    % Comando para añadir nota (invisible en el cuerpo)
    \newcommand{\eqnote}[1]{%
      \addcontentsline{leq}{leqnote}{#1}%
    }
    
    % Comando para ecuaciones
    \newcommand{\eqblock}[2]{%
      % Verificar si necesitamos imprimir el título de sección
      \ifx\leq@current@section\leq@last@printed@section\else
        \addcontentsline{leq}{leqsection}{\leq@current@section}%
        \global\let\leq@last@printed@section\leq@current@section
        \gdef\leq@last@printed@subsection{}% Reset subsección al cambiar sección
      \fi
      % Verificar si necesitamos imprimir el título de subsección
      \ifx\leq@current@subsection\leq@last@printed@subsection\else
        \ifx\leq@current@subsection\@empty\else
          \addcontentsline{leq}{leqsubsection}{\leq@current@subsection}%
          \global\let\leq@last@printed@subsection\leq@current@subsection
        \fi
      \fi
      % Ecuación
      \refstepcounter{eqnum}%
      \begin{equation*}
        #1\tag{E.\theeqnum}
      \end{equation*}%
      \addcontentsline{leq}{leqt}{[E.\theeqnum] -- #2}%
      \addcontentsline{leq}{leqm}{#1}%
    }
    
    %%% ----- Índice de ecuaciones ----------%%%
    % Tamaño para []
    \newcommand{\leqIndexFont}{\normalsize}
    \newcommand{\leqTitleFont}{\tiny}
    \newcommand{\leqNoteFont}{\tiny}
    % Cabecera de sección 
    \newcommand*\l@leqsection[2]{%
      \addvspace{25pt}% Mayor separación antes de nueva sección
      \noindent\leqIndexFont
      \makebox[\linewidth]{\rule{.38\linewidth}{0.4pt}\ \ \textbf{  \MakeUppercase{#1}}\ \ \rule{.38\linewidth}{0.4pt}}%
      \par\vspace{-10pt}
    }
    % Cabecera de subsección 
    \newcommand*\l@leqsubsection[2]{%
      \par\addvspace{15pt}%
      \noindent\leqIndexFont #1
      \par\vspace{-7pt}
    }
    % Título de ecuación 
    \newcommand*\l@leqt[2]{%
    \par\addvspace{10pt}%
      \par\noindent\leqTitleFont #1\par
    }
    % Miniatura de ecuación
    \newcommand*\l@leqm[2]{%
      \vspace{0.3\baselineskip}%
      \par\scriptsize\noindent\hspace*{1.5em}\(#1\)\par%
    }
    % Nota de ecuación 
    \newcommand*\l@leqnote[2]{%
      \vspace{0.2\baselineskip}%
      \par\noindent\leqNoteFont\hspace*{1.5em}\textbf{Nota.--} #1\par\vspace{0.5\baselineskip}%
    }
    % Generación del índice
    \newcommand{\mylistofequations}{%
      \clearpage
      \phantomsection
      % Título visual de la lista (ajústalo a tu gusto):
      {\centering\bfseries\Large Índice de ecuaciones\par}
      % Entrada en nuestro índice 'stc':
      \addcontentsline{stc}{section}{Índice de ecuaciones}%
      % Llamada a tu lista real (usa el nombre exacto de tu macro):
      \@starttoc{leq}%
    }
    
    %--------- Captura de títulos de sección --------%
    \let\leq@original@section\section
    \let\leq@original@subsection\subsection
    % Flag para saber si estamos en una sección asterisco
    \newif\ifLeqStarSection
    \LeqStarSectionfalse
    
    % Redefinir \section CON CONTROL DE ASTERISCO
    \renewcommand{\section}{%
      \@ifstar{\leq@section@star}{\@dblarg\leq@section@nostar}%
    }
    \def\leq@section@star#1{%
      \global\LeqStarSectiontrue
      \gdef\leq@current@section{#1}%
      \gdef\leq@current@subsection{}%
      \leq@original@section*{#1}%
      \global\LeqStarSectionfalse
    }
    \def\leq@section@nostar[#1]#2{%
      \global\LeqStarSectionfalse
      \gdef\leq@current@section{#2}%
      \gdef\leq@current@subsection{}%
      \leq@original@section[#1]{#2}%
    }
    % Redefinir \subsection
    \renewcommand{\subsection}{%
      \@ifstar{\leq@subsection@star}{\@dblarg\leq@subsection@nostar}%
    }
    \def\leq@subsection@star#1{%
      \global\LeqStarSectiontrue
      \gdef\leq@current@subsection{#1}%
      \leq@original@subsection*{#1}%
      \global\LeqStarSectionfalse
    }
    \def\leq@subsection@nostar[#1]#2{%
      \global\LeqStarSectionfalse
      \gdef\leq@current@subsection{#2}%
      \leq@original@subsection[#1]{#2}%
    }

    % ----- Restauración de valores Globales de estilo ----------%
    % Flags
    \newif\ifInFootnote
    \InFootnotefalse
    \pretocmd{\@footnotetext}{\InFootnotetrue}{}{}
    \apptocmd{\@footnotetext}{\InFootnotefalse}{}{}
    % Profundidad de listas (soporta anidadas)
    \newcount\MyListDepth \MyListDepth=0
    \AtBeginEnvironment{la}{\advance\MyListDepth by 1}
    \AfterEndEnvironment{la}{\advance\MyListDepth by -1}
    \AtBeginEnvironment{lnum}{\advance\MyListDepth by 1}
    \AfterEndEnvironment{lnum}{\advance\MyListDepth by -1}
    \AtBeginEnvironment{itemize}{\advance\MyListDepth by 1}
    \AfterEndEnvironment{itemize}{\advance\MyListDepth by -1}
    \AtBeginEnvironment{enumerate}{\advance\MyListDepth by 1}
    \AfterEndEnvironment{enumerate}{\advance\MyListDepth by -1}
    % Restauración diferida: hazlo al INICIO del próximo párrafo real
    \newif\if@pendingRestore
    \@pendingRestorefalse
    \newcommand{\RequestRestore}{\global\@pendingRestoretrue}
    \everypar{%
      \if@pendingRestore
        \global\@pendingRestorefalse
        \ifInFootnote\else
          \ifnum\MyListDepth=0
            \setlength{\parskip}{\GlobalParskip}%
            \setstretch{\GlobalBaseLineStretch}%
          \fi
        \fi
      \fi
    }
    % Al final de bloques:
    \AfterEndEnvironment{equation}{\RequestRestore}
    \AfterEndEnvironment{equation*}{\RequestRestore}
    \AfterEndEnvironment{la}{\RequestRestore}
    \AfterEndEnvironment{lnum}{\RequestRestore}
    % Al final de macros:
    \apptocmd{\eqblock}{\RequestRestore}{}{}
    \apptocmd{\imagenfit}{\RequestRestore}{}{}
    
\makeatother

% ===================== ToC propio con 4 niveles (único bloque) =====================
\makeatletter

% --- Escritores: añadimos una línea al .stc en cada cabecera numerada ---
% Guardamos las versiones actuales para llamar al comportamiento normal
\let\MyTOC@orig@section\section
\let\MyTOC@orig@subsection\subsection
\let\MyTOC@orig@subsubsection\subsubsection
\let\MyTOC@orig@paragraph\paragraph

% SECTION
\renewcommand{\section}{%
  \@ifstar{\MyTOC@section@star}{\@dblarg\MyTOC@section@nostar}%
}
\def\MyTOC@section@star#1{%
  \MyTOC@orig@section*{#1}% sin entrada al .stc
}
\def\MyTOC@section@nostar[#1]#2{%
  \MyTOC@orig@section[{#1}]{#2}% comportamiento normal (escribe al .toc)
  % Escribimos también nuestra línea al .stc (texto con \numberline)
  \addcontentsline{stc}{section}{\protect\numberline{\thesection}#2}%
}

% SUBSECTION
\renewcommand{\subsection}{%
  \@ifstar{\MyTOC@subsection@star}{\@dblarg\MyTOC@subsection@nostar}%
}
\def\MyTOC@subsection@star#1{%
  \MyTOC@orig@subsection*{#1}%
}
\def\MyTOC@subsection@nostar[#1]#2{%
  \MyTOC@orig@subsection[{#1}]{#2}%
  \addcontentsline{stc}{subsection}{\protect\numberline{\thesubsection}#2}%
}

% SUBSUBSECTION
\renewcommand{\subsubsection}{%
  \@ifstar{\MyTOC@subsubsection@star}{\@dblarg\MyTOC@subsubsection@nostar}%
}
\def\MyTOC@subsubsection@star#1{%
  \MyTOC@orig@subsubsection*{#1}%
}
\def\MyTOC@subsubsection@nostar[#1]#2{%
  \MyTOC@orig@subsubsection[{#1}]{#2}%
  \addcontentsline{stc}{subsubsection}{\protect\numberline{\thesubsubsection}#2}%
}

% PARAGRAPH
\renewcommand{\paragraph}{%
  \@ifstar{\MyTOC@paragraph@star}{\@dblarg\MyTOC@paragraph@nostar}%
}
\def\MyTOC@paragraph@star#1{%
  \MyTOC@orig@paragraph*{#1}%
}
\def\MyTOC@paragraph@nostar[#1]#2{%
  \MyTOC@orig@paragraph[{#1}]{#2}%
  % Si no numeras los \paragraph, cambia \theparagraph por texto plano sin \numberline
  \addcontentsline{stc}{paragraph}{\protect\numberline{\theparagraph}#2}%
}

% --- Impresor del índice propio (lee el .stc) ---
\newcommand{\MyTOC}{%
  \par\bigskip
  {\centering\bfseries\Large Índice de contenido\par}%
  \medskip
  \begingroup
    \parindent=0pt\parskip=2pt
    \@starttoc{stc}%
  \endgroup
}

\makeatother
% ===================== fin ToC propio =====================


%%%%%%%%%%%%%%


% --- Datos del tema ---
\title{Tema 3.A.32 -- El modelo de oferta y demanda agregada.\\
\large Determinación de renta e inflación en una economía abierta. Análisis de las políticas monetaria y fiscal, y de los shocks y políticas de oferta}
\author{Marcelo Ortega}
\date{Marzo 2026}
\newcommand{\TemaCode}{3A32}

\oldtitle{Tema 3.A.28. Determinación de renta, precios e inflación en una economía abierta. Análisis de las políticas monetaria y fiscal.

          Tema 3.A.40. \textit{Shocks} de oferta: modelos e implicaciones de política económica.}
\changerationale{
    Este tema tiene por objeto recoger cómo impactan los shocks en un sistema económico en equilibrio agregado (origen, propagación, efectos de segunda ronda), así como cómo pueden interrelacionarse shocks de oferta y de demanda contemporáneos (como ocurrió, por ejemplo, a raíz de la pandemia de COVID-19). El opositor debería conocer qué respuesta de política económica ha de aplicarse en función de las características del shock. Resultaría deseable para el opositor estudiar este tema después de los A.42, y A.36 a A.41. La referencia a una economía abierta ha de entenderse como una extensión, por ser muy frecuentes los shocks de oferta y de demanda procedentes del resto del mundo. El grueso del contenido del anterior tema A.28 se desplaza al bloque B, al tratarse de cuestiones de Macroeconomía abierta.
}

\usepackage{soul}\sethlcolor{yellow}% resaltado amarillo: \hl{texto} (Panel TCEE, ⌃H)
\begin{document}


\maketitle
\changebox

\newpage

% --- Página de resumen y modificaciones ---
\puntosclave{ %% Escribir puntos clave Apuntes CECO
    }
\vspace{1cm}

\modificaciones{
    
    }

\newpage
\MyTOC
\newpage

\mylistofequations
\clearpage

\MyListOfFigures
\clearpage


\pagenumbering{arabic}
\setcounter{page}{1}


\section{Introducción}\label{intro}


\subsection{Contextualización}\label{context}


\subsubsection{Relevancia}\label{importance}


\subsubsection{Historia}\label{history}



\subsubsection{Problemática}\label{problem} 




\subsection{Estructura de la exposición}\label{structure}
La sección 2 revisa los instrumentos para el análisis empírico de las perturbaciones exógenas y los principales resultados de su aplicación para la identificación de shocks de oferta y de demanda. La sección 3 detalla el modelo básico neokeynesiano que se utilizará para el análisis de los efectos de los diversos shocks. Las secciones 4 y 5 se centran en las predicciones del modelo teórico ante estas perturbaciones, su consistencia con los datos y sus implicaciones de política económica. En la sección 6 se presenta una extensión al caso de una economía abierta. El capítulo termina con unas conclusiones.



\clearpage
\section{Modelo canónica de la NOEM: Determinación de la renta y la inflación en una economía abierta}
\puntosclave{
    El modelo neo-keynesiano es el principal marco de análisis macroeconómico actual. 
    
    Es un modelo de equilibrio general microfundamentado con rigideces nominales compuesto por hogares, empresas y gobierno o banco central. 
    
    El modelo en su versión más sencilla se reduce a tres ecuaciones: curva IS dinámica, la nueva curva de Phillips neo-keynesiana y la regla de política monetaria. La introducción de supuestos simplificadores frecuentemente empleados nos pem1ite una representación del modelo a través de dos ecuaciones, muy útil para modelizar los shocks en función de si su procedencia es de oferta o de demanda.
}   

Gran parte del trabajo reciente en macroeconomía se ha centrado en el desarrollo y evaluación de modelos monetarios que incorporan \textbf{competencia imperfecta y rigideces nominales} dentro de una estructura de equilibrio general dinámico estocástico (DSGE), el modelo \textit{workhorse} de la teoría de los ciclos reales (RBC). 

El resultado de este sincretismo ha sido incorporar los efectos no triviales que tienen los cambios en las condiciones monetarias sobre las variables reales. La Política Monetaria puede así convertirse en una herramienta potencial de estabilización, así como en una fuente independiente de fluctuaciones económicas. No resulta sorprendente, por tanto, que el estudio de las propiedades de distintas reglas de política monetaria haya sido un área de investigación fructífera en los últimos años y una aplicación natural de esta nueva generación de modelos.

A lo largo de la exposición nos apoyearemos en uno de los modelos más populares de la NOEM: un modelo dinámico, estocástico de equilibrio general para una economía pequeña y abierta propuesto por GALÍ y MONACELLI. 

\subsection{Economía pequeña y abierta: GALÍ y MONACELLI (2005)}
El Modelo de GALÍ y MONACELLI parte de una estructura micro-fundamentada con agentes que optimizan sus decisiones dada la información de que disponen en cada momento del tiempo $t$. Podemos agrupar estos agentes en \textbf{hogares, empresas, gobierno y banco central}. Los hogares en este modelo consumen, trabajan y ahorran en un conjunto homogeneo de activos que rinden a un determinado tipo de interés real. La decisión óptima consumo ahorro surge de la maximización del flujo de utilidades presente y futuras sujeta a una restricción presupuestaria.\footnote{%
    Por motivos de tiempo, en la exposición no conviene derivar analíticamente las ecuaciones, simplemente mostrar sus ecuaciones principales y el desarrollo gráfico. Puede consultarse la derivación analítica en [\authorfont{Ver Anexo \ref{anx:nek2_gali}}].
}

El equilibrio se deriva de la interacción entre dos ecuaciones:

\eqblock{\begin{aligned}
    &\text{Curva IS}:&&\boxed{x_t = \mathbb E_t[x_{t+1}] - \frac{1}{\sigma_\alpha}\big(i_t - \mathbb E_t[\pi_{H,t+1}] - \overline{rr}_t\big)}\\[2em]
    &\text{Curva de PHILLIPS}:&&\boxed{\pi_{H,t} = \beta \mathbb E_t[\pi_{H,t+1}] + \kappa_\alpha x_t}
\end{aligned}
}{Modelo canónico de la NOEM. GALI y MONACELLI (2005)}

Por tanto, la dinámica hacia el equilibrio surge de la interacción de estas las ecuaciones que muestran la interacción entre las principales variables macroeconómicas de una economía pequeña y abierta, en representación incremental (o log-linearizada). Como vemos, puede representarse en términos de brecha de producción ($x_t$) e inflación doméstica ($\pi_{H,t}$), una representación análoga a la que los modelos NEK-2 clásicos ofrecen para una economía cerrada.

\subsubsection{Curva de PHILLIPS: Optimización de las empresas}
La forma de la ecuación de PHILLIPS en la economía abierta es análoga a la de la economía cerrada en lo que respecta a la inflación doméstica. No obstante, existen algunas particularidades características:

\eqblock{
\begin{aligned}
    &\boxed{\kappa_\alpha\equiv\lambda(\sigma_\alpha+\varphi)}\qquad \text{donde,}
    \left\{\begin{aligned}
        &\varphi\equiv\text{Impacto sobre el empleo ($n_t$)}\\[0.5em]
        &\sigma_\alpha(\alpha,\eta)\equiv\text{Términos de intercambio}\\[0.5em]
        &\lambda=\frac{(1-\beta\theta)(1-\theta)}{\theta}\equiv\text{Grado de rígideces locales}
    \end{aligned}\right.\\[1em]
    &\text{De hecho, si}\quad\underset{\alpha\in[0,1]}{\alpha\to0}\Longrightarrow\kappa\equiv\lambda(\sigma+\varphi)\sim\text{NEK-2}\\[1em]
    &\left[
    \begin{aligned}
        &\text{\scriptsize $\underset{\text{\scriptsize OJO, sin entrar a valorar economía abierta o cerrdad: solo perspectiva del consumidor}}{\text{Es útil recordar que la CPO del problema del consumidor es:}}$}\\[0.5em]
        &\text{\scriptsize $C_t^\alpha N_t^\varphi=\frac{W_t}{P_t}\Longrightarrow w_t-p_t=\sigma c_t+\varphi n_t\Longrightarrow c_t=\mathbb E_t[c_{t+1}]-\frac{1}{\sigma}(i_t-\mathbb E_t[\pi_{t+1}]-\rho)$}
    \end{aligned}\right]\\[1em]
    &\text{¿Por qué?}:\quad \sigma_\alpha=\frac{\sigma}{1-\alpha+\alpha\omega}\Longrightarrow
    \left\{\begin{aligned}
        &\alpha\to0:\frac{\partial \pi_{H,t}}{\partial x_t}\to\sigma\\[1em]
        &\sigma_\alpha<\sigma,\quad \forall \alpha\in[0,1]
    \end{aligned}\right.
\end{aligned}
}{}

Por tanto, como vemos el grado de apertura ($\alpha$) afecta a la dinámica de la inflación a través de su impacto sobre la pendiente de la curva de PHILLIPS, es decir, sobre la sensibilidad de la inflación ante variaciones en la brecha de producción. ¿Por qué?
\begin{lnum}
    \item \textbf{Empleo}. En primer lugar porque en una economía abierta alteraciones en la brecha de producción se transmiten a través del empleo a los costes marginales (la elasticidad del empleo se captura por $\varphi$); y,
    \item \textbf{Términos de intercambio}. En segundo lugar, porque afecta a los costes marginales a través de los términos de intercambio (capturados por $\sigma_\alpha$). En particular, si $\sigma\eta > 1$,\footnote{%
        Cada uno de estos parámetros corresponde a  grado de sustituibilidad intertemporal del consumo ($\sigma$) y grado de sustituibilidad entre (producto extranjero)v vs. (producto nacional), $\eta>0$)
    } una mayor apertura reduce el ajuste necesario de los términos de intercambio para absorber alteraciones en la brecha de producción (relativo al mundial), amortiguando así su impacto sobre el coste marginal y la inflación.
\end{lnum}

\textcolor{red}{\textbf{OJO!} Es importante enfatizar que uno de las diferencias más importantes entre la especificación del modelo para una economía cerrada y para una economía abierta se da precisamente en la especificación de los costes marginales, que origina la interrelación entre output gap e inflación (costes marginales es proporcional a output gap). Esto es así porque, en economía abierta, los costes marginales continen el parámetro $-\nu$. No me ha quedado claro el impacto que tiene esta diferencia por lo que se recomienda mirar el modelo de FURLANETTO y GALÍ y MODIGLIANI para encontrar esto. Está estrechamente relacionado con el subsidio al desempleo que los autores especifican, $-\nu=-\log(1+\tau)$}

En definitiva, $\sigma_\alpha$ desempeña el papel de inversa de la elasticidad de sustitución intertemporal, teniendo en cuenta que es posible sustituir el consumo no solo intertemporalmente, sino también entre países. Además, $\sigma_\alpha$ es siempre menor que $\sigma$ y es igual a $\sigma$ en el caso de una economía cerrada.

\subsubsection{Curva IS: Optimización de los Hogares}
Respecto a la Curva IS, se expresa en términos de la brecha de producción, igual que en el canónico de la NEK-2 y de tipo forward looking. No obstante, hay dos diferencias importantes:

\eqblock{
\begin{aligned}
    &\text{Tipo de interés natural}:\quad \overline{rr}_t(\text{NOEM})\neq r_t^n(\text{NEK})\\[1em]
    &\text{En especial},
    \left\{\begin{aligned}
        &\Delta\alpha\Rightarrow \underset{\text{\scriptsize sensitividad}}{|\Delta Cov(x_t,i_t)|}{\iff} \omega>1\quad(\Leftarrow\underset{\omega(\eta,\gamma)}{\Delta\eta/\Delta\gamma)}\\[1em]
        &\Delta\alpha\Rightarrow {|\Delta Cov(\overline{rr}_t,\mathbb E_t[y_t^*])|}
    \end{aligned}\right.\\
    &\left[\text{\scriptsize $\overline{rr}_t=\rho-\sigma_\alpha\Gamma(1-\rho_\alpha)a_t+\alpha\sigma_\alpha(\Theta+\Psi)\mathbb E_t[\Delta y_{t+1}^*]$}\right]
\end{aligned}
}

Como vemos, la diferencia fundamental se encuentra: (i) en la determinación del tipo de interés natural de dicha economía; y, (ii) a cómo afecta la apertura a la sensibilidad (pendiente) ante variaciones de la brecha entre tipo real y tipo natural. ¿Por qué?
\begin{lnum}
    \item \textbf{Tipo de interés}. En primer lugar, el grado de apertura altera la sensibilidad de la brecha de producción ante variaciones en el diferencial de tipos de interés. En particular, si $\omega > 1$ (que se obtiene con valores \textit{altos} de $\eta$ y $\gamma$), una mayor apertura aumenta dicha sensibilidad (recordar que $\sigma_\alpha<\sigma$), debido al papel amplificador de los términos de intercambio sobre la demanda. Es decir, ;\footnote{%
        $\gamma$ es el grado de sustituibilidad entre productos producidos en diferentes economías extranjeras. Por lo tanto, vemos que es sustancialmente diferente al parámetro $\eta$ aunque es una consecuencia directa de la apertura.
    }
    \item \textbf{Interrelaciones comerciales}. En segundo lugar, la apertura introduce dependencia del tipo de interés natural respecto al crecimiento esperado del output mundial, además de la productividad doméstica.
\end{lnum}

\subsubsection{Características básicas: Parámetros, Supuestos y Política Monetaria}
En definitiva, el modelo sobre el que nos apoyaremos es un tipo DSGE para una economía abierta pequeña extremadamente detallado. En esencia, incorpora los siguientes supuestos:
\begin{lnum}
    \item \textbf{Rigideces nominales}. Introducir rigideces de precios en la forma popularizada por CALVO (1983), como hacen GALI y MONACELLI (2005), nos permite estudiar los efectos de un shock persistente en un marco coherente con las funciones de respuesta al impulso derivadas de la evidencia VAR;
    \item \textbf{Tipo de cambio real y nominal dinámicos}. Como vemos, la introducción de un ligero sesgo hacia el consumo doméstico permite desviaciones respecto a la PPP (que se cumple en Estado Estacionario).\footnote{%
        Este comportamiento es casi una elección natural en un marco de economía pequeña y abierta: de lo contrario, en caso de preferencias idénticas entre países, dado que la participación del bien doméstico en la cesta de consumo de los países extranjeros es irrelevante, el país doméstico debería consumir solo bienes extranjeros y exportar toda la producción del bien doméstico.
    } Las desviaciones respecto a la PPP, junto con la UIP (ahora veremos), permiten una dinámica rica en la respuesta del tipo de cambio real y nominal ante perturbaciones.
    \item \textbf{Mercados financieros completos}. La existencia de mercados financieros completos permite que, en Estado Estacionario (equilibrio), se cumpla la Paridad No Cubierta de Intereses (PNCI) a la vez que permite dinámicas de ajuste (respuestas de impulso) que flexibilizan la dinámica del Tipo de Cambio real y nominal;
    \item \textbf{Preferencias y tecnología neoclásicas de buen comportamiento}; y,
    \item \textbf{Política Monetaria}. El supuesto de una reacción \textbf{endógena} de la Política Monetaria implica una respuesta muy diferente de la economía ante los shocks en función del tipo de regla a la que se adhiera. Por tanto, la asuencia explícita de esta permite precisamente evaluar los diferentes regímenes y encontrar alguno que sea óptimo, siendo, en todo caso, el tipo de interés a corto plazo el instrumento principal.
\end{lnum}

Naturalmente los coeficientes en la economía abierta dependen también de parámetros específicos (como el grado de apertura y la sustituibilidad entre bienes de distinto origen), mientras que las perturbaciones incluyen fluctuaciones del output mundial (tomadas como exógenas).

\paragraph{Implicaciones de Política Económica: \textit{Inflation targeting}}
Cabe hacer especial énfasis en que los resultados de estos autores muestran que estos regímenes pueden ordenarse según la \textbf{volatilidad del tipo de cambio nominal} y de los \textbf{términos de intercambio} que generan. 

Los autores concluyen que, en términos de bienestar, la Política Monetaria óptima consiste en mantener un \textbf{objetivo estricto de inflación doméstica}.\footnote{%
    En términos de Política Económica, esto nos permite analizar las implicaciones macroeconómicas de diferentes regímenes de Política de Monetaria basados en reglas para la pequeña economía abierta. En particular, los autores analizan tres: reglas de Taylor basadas en la inflación doméstica y en el IPC, y un régimen de tipo de cambio fijo.
} 

Es por esto que a lo largo de nuestra exposición consideraremos la Autoridad Monetaria mantiene esta regla.\footnote{%
    Los autores demuestran que una diferencia clave entre estos regímenes radica en el grado de volatilidad del tipo de cambio que implican. La excesiva suavidad del tipo de cambio bajo el resto de reglas (en comparación con \textit{inflation targeting}), junto con la rigidez de precios nominales, impide que los precios relativos se ajusten suficientemente rápido ante shocks de productividad relativos, generando desviaciones importantes respecto al óptimo de primer orden. No obstante, una regla de Taylor basada en el IPC da lugar a una dinámica intermedia entre una regla basada en inflación doméstica y un tipo de cambio fijo.

Por otro lado, la volatilidad de los términos de intercambio se traduce directamente en un ordenamiento en términos de bienestar. Para un amplio rango de parámetros, una regla de Taylor basada en inflación doméstica domina a una basada en el IPC, y esta a su vez domina a un tipo de cambio fijo. De forma general, muestran que cuanto mayor es la volatilidad de los términos de intercambio en equilibrio, menor es la volatilidad de la inflación y de la brecha de producción, y mayor es el bienestar.

En definitiva, la política óptima implica estabilizar completamente el nivel de precios domésticos.
} No obstante, cabe recordar que una política de objetivo estricto de inflación doméstica implica una volatilidad mucho mayor del tipo de cambio y de los términos de intercambio que la que se obtiene bajo las reglas de Taylor o un Tipo de Cambio fijo. 

\subsection{Extensiones: Instrumentos y Política Fiscal}
Por último, antes de pasar a hablar de los \textit{shocks} que afectan a una economía abierta y cuáles son las herramientas de Política Económica de las que disponemos cuando debemos reacciones, debemos introducir dos elementos principales. 

\subsubsection{Análisis gráfico: Aspa de la NOEM}
En primer lugar, el sistema (o modelo) que acabamos de formular se puede estudiar como recoge el gráfico siguiente:

\imagenfit{Grafico_NOEM.png}{Representación gráfica del modelo NOEM: GALÍ y MONACELLI (2005)}{Apuntes ICEX-CECO. Tema 3.A.32}
  
El origen de coordenadas (intersección) marca el equilibrio en el modelo una vez calibrados los parámetros fundamentales de nuestra economía (y mundo). De esta forma podremos analizar de forma muy intuitiva cómo afectan los shocks de diferente naturaleza y cómo se transmiten en términos dinámicos. En esencia, lo que estamos representado es una suerte de Estado Estacionario del que la economía no se desviaría sino cambiarán sus fundamentales o se viera afectada por shocks que afectan tanto a los componentes como a las previsiones de los agentes. 

\subsubsection{Política Fiscal: FURLANETTO (2006)}
Por otro lado, hay algo de lo que carece sonadamente el modelo anteriormente expuesto: la Política Fiscal. Aunque es cierto que los autores introducen como parte de la restricción presupuestaria del consumidor la existencia de impuestos y subvenciones (aunque configuradas como no distorsionantes) y, además, introducen una subvención/apoyo al empleo que reduce los costes marginales de los productores (neutralizando así su poder de mercado).\footnote{%
    En conreto, se introduce el supuesto de un subsidio al empleo constante $\tau$ ($\nu=\log{(1-\tau)}$) que neutraliza la distorsión asociada al poder de mercado de las empresas.

La intuición de este resultado es sencilla: con el subsidio en vigor, solo queda una distorsión efectiva en la economía, a saber, la rigidez de precios. Al estabilizar los márgenes en su nivel sin fricciones, las rigideces nominales dejan de ser vinculantes, ya que las empresas no tienen incentivos para ajustar sus precios. Por construcción, la asignación de equilibrio resultante es eficiente y el nivel de precios permanece constante.
} 

No obstante, FURLANETTO extendió este modelo un año después para analizar qué ocurría en la economía cuando se producía un shock fiscal, prestando atención a los efectos que esto tendría sobre el output gap, la inflación y el tipo de cambio real. Para ello, introduce un sector público que realizará gasto público asumiendo tres cuestiones clave:
\begin{lnum}
    \item \textbf{Gasto Público}. El gasto público se modeliza como puro gasto: no se pretende analizar los efectos de la Política Fiscal sobre el bienestar, sino ver el efecto que cambios en el tono de la Política Fiscal tienen sobre la dinámica endógena de la economía;
    \item \textbf{Sesgo doméstico} ($\eta_G=0$). El gasto público está totalmente sesgado a la producción doméstica, de forma que el sector público solo compra bienes nacionales; y,
    \item \textbf{Presupuesto equilibrado}. El aumento del gasto público debe compañarse de un aumento equivalente en impuestos. Es decir, el presupuesto está equilibrado en cada período. (Nótese que esto afecta indirectamente al bienestar del consumidor, pero por simplicidad no será tratado) 
\end{lnum}

Por tanto, el otuput estará formado no solo por consumo y exportaciones netas, sino también por el Gasto Público. Las implicaciones analíticas (y gráficas) del Gasto Público son fácilmente deducibles:

\eqblock{
\begin{aligned}
    &(1)\;\;\text{Si},\;\;\frac{g_t}{y_t}\longrightarrow1\quad \Longrightarrow\quad 
    \left\{\begin{aligned}
        &\underset{\uparrow y_t\Rightarrow\uparrow g_t,\;\text{\scriptsize pero, $\eta_G=0$}\Rightarrow \uparrow\pi_{H,t}}{\uparrow Cov(x_t,\pi_{H,t})}\\[1em]
        &\downarrow|Cov(x_t,r_t-\overline{rr}_t)|
    \end{aligned}\right.\\[2em]
    &(2)\;\;Cov(g_t,\overline{rr}_t)>0\quad \underset{\left\{\begin{aligned}
        &\text{\scriptsize efecto riqueza}\\
        &\text{\scriptsize persistencia  $\rho_g$}
    \end{aligned}\right.}{\overset{\substack{\text{\scriptsize a través de:}}}{\Longrightarrow}}\quad \text{Si},\;\;\uparrow g_t \Longrightarrow\quad \uparrow \overline{rr}_t\\
\end{aligned}
}{Sensibilidad paramétrica: Gasto Público. Modelo NOEM: FURLANETTO}

El autor analiza también las dinámicas ante diferentes reglas monetarias pero para simplificar nuestra exposición y mantener la coherencia, expondremos únicamente las implicaciones de mantener una regla \textbf{inflation targeting}.\footnote{%
    Ahora bien, en este caso las implicacinoes de una elección u otro si que tienen significativas diferencias enel impacto de un aumento del Gasto Público. Pueden verse los resultados en el paper original.
} Aunque hay que ser plenamente consciente que al establecer dicho objetivo, las variaciones en el tono de la Política Fiscal tienen un efecto mucho menor que ante una política de tipo de cambio fijo. Esto es así porque, bajo Tipo de Cambio Fijo la Autoridad no elevaría los tipos de interés para frenar las tendencias inflacionarias domésticas y la apreciación real sería menor, por lo que el impacto que tendría sobre el consumo y la balanza comercial se limitaría (por ejemplo, PF expansiva produce una caída del consumo y un empeoramiento de la BP). Lo que, de hecho, confirma el conocido resultado del Modelo MUNDELL-FLEMMING: la Política Fiscal es más efectiva bajo Tipo de Cambio Fijo que bajo un régimen flexible.

\clearpage
\section{Dinámicas de ajuste: Shocks}
\puntosclave{
    
}

Ahora que conocemos el \textbf{caballo de batalla} de la NOEM podemos presentar las implicaciones de Política Económica que esta vertiente proporciona a la hora de reaccionar a shocks de diferente naturaleza. 

\textbf{¿Qué queremos acabar diciendo?} Las perturbaciones (o shocks) que experimentan las economías exigen (en muchas ocasiones) una respuesta por parte del policy maker, y encontrar dicha respuesta óptima no suele ser fácil. El Modelo de la NOEM es tremendamente útil dada su simplicidad y la capacidad que tiene para proporcionar dinámicas de ajuste estables y concretas de los diferentes parámetros (muchos de ellos a disposición del policy maker). No obstante, todos estos resultados deben tomarse con muchísima cautela. ¿Por qué? Principalmente, porque la modelización parámetrica, precisamente por su simplicidad, exige que el investigador descarte (o rechaze) infinidad de parámetros que son (o pueden ser) de absoluta importancia en el análisis económico.

Antes de nada, es importante definir qué se entiende por perturbaciones. La literatura económica tiende a distinguir los shocks atendiendo a dos grandes grupos: 
\begin{la}
    \item \textbf{Perturbaciones de demanda}. Partiendo de la definición de PIB por el lado del gasto, podemos hablar de perturbaciones de demanda originadas en la demanda de bienes de consumo privado, demanda de bienes de inversión y demanda de gasto público; y,
    \item \textbf{Perturbaciones de oferta}. Por otro lado, considerarnos perturbaciones de oferta cambios exógenos en la productividad o tecnología de las empresas y en los costes de producción.
\end{la}

No obstante, debería haber quedado claro mientras desarrollábamos el modelo que la alta sensibilidad que los parámetros mantienen entre sí hace que, a pesar de que el origen de la perturbación pueda ser claramente identificado entre estos, sus efectos son mucho más generales, produciéndose efectos de primera, segunda, tercera ronda, \textit{in fine} (quizás en algunos casos), que no se restringen a los ámbitos que han motivado su aparición. 

\subsection{Perturbaciones de demanda}
Trataremos en primer lugar algunos ejemplos de perturbaciones de demanda. En particular, nos centramos en dos tipos de perturbaciones de demanda que desde la segunda mitad de los años 80's constituyen el centro del análisis de política económica actual y que contribuyen a explicar la mayor parte de las fluctuaciones del PIB en los países industrializados. 

\subsubsection{Shock (positivo): Gasto Público}
¿Qué sucede ante un cambio exógeno del Gasto Público? Por ejempplo, ¿cómo afectaría al equilibrio un incremento del 1\%, temporal y anticipado, del Gasto Público?\footnote{%
    Recordemos que según la evidencia empírica, las perturbaciones de demanda tienen un efecto transitorio en el nivel de producto (BLANCHARD y QUAH, 1989) y prevén una correlación positiva entre producto y precios (GALÍ, 1992).
} 

Sabemos que esta perturbaciones afectaría directamente a la curva IS dinámica del modelo. Un aumento de $g_t$ hoy para volver a su nivel inicial pasado un período. Gráficamente, la Curva IS se desplazaría como sigue:

\imagenfit{PF_Expansivo.png}{Efectos de un aumento exógeno del Gasto Público en el modelo de la NOEM: FURLANETTO}{ Apuntes ICEX-CECO. Tema 3.A.32}

El aumento del gasto público incrementa la demanda total de bienes (domésticos). Esto tiene implicaciones directas tanto desde una óptica nacional como exterior: (Seguimos la explicación de FURLANATTO)
\begin{lnum}
    \item \textbf{Tipo de interés natural}. Se experimenta un incremento del tipo de interés natural de la economía. Esto tiene varias causas, dentre las cuales destacan los efectos generados sobre las decisiones intertemporales de consumo de los agentes;
    \item \textbf{Sector exterior}. Desde un piunto de vista externo, tanto el tipo de cambio efectivo real como el tipo de cambio nominal experimentan una notable apreciación, lo que empeora notablemente los términos de intercambio debido, principalmente, al sesgo doméstico que sufre el Gasto Público. Además, esto afecta al deterioro de la Balanza Comercial a través del canal consumo (sustitución por importaciones);
    \item \textbf{Consumo}. El consumo experimenta una caída debido a: (i) el deterioro de sus ingresos debido a la necesidad de mantener un presupuesto equilibrado; y, (ii) la sustitución intertemporal de consumo debido al incremento de los tipos de interés como reacción de la Autoridad Monetaria. No obstante, esta caída es muy ligera y destaca el hecho de que a mayor grado de apertura menor será la caída del consumo puesto que el consumo doméstico será sustituido por consumo en bienes importados (lo que, en consecuencia, explica el deterioro de la Balanza Comercial);
    \item \textbf{Producción}. Por lo que respecta al output, el multiplicador de la Política Fiscal resulta ser sustancialmente menor a 1 (en torno al $0,5$). Por lo tanto, si bien se xperimentaría un incremento del output inmediato, no sería tan pronunciada como cabría esperar atendiendo a la Teoría Keynesiana;\footnote{%
        Surgen dos resultados importantes:
    \begin{lnum}
        \item El multiplicador del nivel de producción es en gran medida inferior a uno, tal como encuentra PEROTTI (2004); y,
        \item Dicho multiplicador \textbf{disminuye con el grado de apertura}. En el impacto inicial, nuestras funciones impulso-respuesta estiman un multiplicador de nivel de aproximadamente $0,45$ en el caso de economía cerrada y de alrededor de $0,2$ en la economía abierta bajo nuestra calibración base.
    \end{lnum}
    Estos resultados son consistentes con un efecto negativo significativo sobre las exportaciones netas, como se muestra en la figura 1b y como documentan Giuliodori y Beetsma (2004) y Lane y Perotti.
    }
    \item \textbf{Inflación}. En principio, el shock fiscal no tiene efecto sobre la inflación esperada doméstica porque el banco central aumenta el tipo de interés alcanzando el objetivo de inflación doméstica cero (no obstante, la inflación no esperada queda indeterminada y cabe esperar una ligera subida, ceteris paribus). De hecho, cabría esperar que el efecto sobre el IPC fuese incluso negativo, ya que los precios de importación disminuyen como consecuencia de la apreciación del tipo de cambio.
\end{lnum}

Todo se puede observar en los siguientes gráficos de respuesta, donde vemos claramente que la perturbación tiene una duración media de $3$ años sobre las variables. 

\imagenfit{ShockFiscal_Regresiones.png}{Respuesta de inercia ante un aumento del Gasto Público}{FURLANETTO (2006)}

\paragraph{Implicaciones de Política Económica}
Las implicaciones en términos de Política Económica son entonces claros: se confirman los principales resultados de la literatura de MUNDELL-FLEMING y, en particular, se espera una apreciación del tipo de cambio nominal y real.\footnote{Cabría esperar un mayor efecto expansivo bajo tipos de cambio fijos.
} Sin embargo, no parece ser tan favorable respecto a la Hipótesis de la Equivalencia Ricardiana, sosteniendo que la bajada de consumo no sería tan pronunciada por ésta debido a la capacidad de sustituir consumo entre bienes domésticos y bienes extranjeros.

La evidencia empírica parece contrastar la validez de estas predicciones, que reproducen un \textbf{multiplicador} del nivel de producción bajo (inferior a uno) y \textbf{decreciente} con el grado de \textbf{apertura} de la economía, consistente con una reacción negativa de las exportaciones netas.

\subsubsection{Shock (negativo): Perturbaciones de preferencias}
Otro tipo de perturbaciones de demanda se puede encontrar en cambios exógenos en la utilidad marginal del consumo en uno o varios períodos $t$, o en cambios en el factor de descuento intertemporal ($\rho$). Este segundo tipo de shocks incluiría por tanto, por ejemplo, un aumento de la paciencia (impaciencia) de los consumidores que les lleva a demandar menos (más) hoy y, por tanto, reduciendo su consumo (ahorro). Analicemos un shock negativo de preferencias temporal no anticipado. 

En este caso, la curva IS se desplaza hacia abajo a la izquierda. La brecha de producción disminuye así como el ivel de precios (inflación). De nuevo, ambas variables muestran una correlación positiva ante una perturbación de demanda.

\textbf{INCLUIR GRÁFICO}

El principal mecanismo de transmisión en este tipo de shocks es la caída del tipo de interés natural de la economía ($\overline{rr}_t$), aunque evidentemente no es el único parámetro afectado y, debemos recordar, que $\sigma_\alpha$ intensifica la desviación. Recordemos que:

\eqblock{
\begin{aligned}
    &\downarrow \rho\Longrightarrow\downarrow\overline{rr}_t:&&\underset{\text{\scriptsize aumenta la brecha entre tipo real efectivo y tipo natural}}{(i_t-\mathbb E_t[\pi_{H,t}]-\overline{rr}_t)>0}\\[1em]
    &\downarrow\overline{rr}_t\Longrightarrow\downarrow x_t\Longrightarrow\downarrow\pi_{H,t}:&&\text{¿Qué hacer?}\\[0.5em]
    &\text{AM}:\quad\downarrow i_t\Longrightarrow\;(i_t-\mathbb E_t[\pi_{H,t}]-\overline{rr}_t)=0,&& \text{pero, si ELB}:\quad \not\downarrow\;i_t\\[0.5em]
\end{aligned}
}{}

Si, en cambio, la autoridad monetaria sigue de cerca el tipo natural, entonces el shock puede absorberse en gran parte via tipo de cambio, términos de intercambio y cuenta corriente, con poca variacion de inflación doméstica y output gap. Ahora bien, si esto no puede realizarse, entonces podemos analizar la serie de parámetros que se verían comprometidos teniendo siempre presente que estos se acentuarán cuanto mayor grado de apertura, pues determina la pendiente de la curva IS:
\begin{lnum}
    \item \textbf{Demanda interna}. Como hemos dicho, un cambio en las preferencias de los agentes afectará directamente al tipo de interés natural de la economía y, en consecuencia, provocará una caída de la demanda interna que deberá ser corregida por la Autoridad Monetaria;
    \item \textbf{Demanda externa}. Desde un piunto de vista externo, el diferencial de tipos reales induce una depreciación de cara al exterior. Es decir, un movimiento positivo de la balanza comercial y un movimiento positivo de los términos reales de intercambio (si $\omega>1$). En una economia mas abierta ese amortiguador externo es mas fuerte, porque $\partial nx_t/\partial s_t=(1-\gamma)\alpha(\frac{\omega}{\sigma}-1)$ aumenta con $\alpha$ y con $\omega$. 
\end{lnum}

Luego hay dos fuerzas simultaneas: por un lado, mayor apertura hace mas sensible la demanda agregada al tipo real via $1/\sigma_\alpha$; por otro, tambien hace mas potente el canal de sustitucion externa via $nx_t$. Si la autoridad monetaria no recorta suficientemente el tipo nominal ($i_t$), cae la brecha de produccion ($x_t$) y cae la inflacion domestica ($\pi_{H,t}$); la depreciacion asociada al diferencial de tipos reales mejora las exportaciones netas ($nx_t(s_t)$) y amortigua la contraccion; \textbf{cuanto mayores son $\alpha$ y $\omega$, mas fuerte es ese canal de \textit{expenditure switching} y mas sensible es la demanda al tipo real}. 

\paragraph{Implicaciones de Política Económica: Gran Recesión y Política Fiscal}
El análisis de estas perturbaciones se emplea en la literatura para explicar la fuerte caída de los tipos de interés durante la Gran Recesión. En este contexto, los agentes se vuelven más pacientes, lo cual reduce el consumo y aumenta fuertemente el ahorro.\footnote{%
    El aumento de los ahorros, en panicular por motivo precaución, también puede darse como respuesta a otro tipo de perturbaciones como por ejemplo un aumento de la aversión al riesgo o un aumento de la incertidumbre.
} 

El aumento de la oferta de ahorro genera una presión a la baja de los tipos de interés nominales. En un contexto de recesión, tales como las que se vivieron entre 2007 y 2008, la autoridad monetaria responderá relajando las condiciones de financiación de los agentes con bajadas adicionales de los tipos de interés. Si esta perturbación es muy persistente, es entonces cuando los Tipos de Interés nominales pueden alcanzar su límite inferior, el conocido por sus siglas en inglés ELB. La experiencia desde la Gran Recesión ha demostrado que en niveles de tipos de interés nominales tan bajos, la Política Monetaria puede dejar de ser efectiva, requiriendo de otros impulsos a la demanda, tales como la Política Fiscal para revertir la situación.\footnote{%
    Más recientemente, en el período post-pandemia, la gran bolsa de- ahorro generada durante el confinamiento, junto con las fuertes políticas tanto fiscal como monetaria de inyecciones de liquidez favorecieron la situación contraria a la de la Gran Recesión: un aumento fuerte de la demanda y de la inflación, especialmente en algunas economías como la estadounidense.
}

\paragraph{Caso particular: Estancamiento secular}
Una aplicación particular y muy relevante de este tipo de perturbaciones es el conocido como \textbf{estancamiento secular}.\footnote{%
    [\authorfont{Ver Tema 3.A.33}] Según la teoría de Keynes, ante un crecimiento de renta sostenido, la propensión media al consumo se reduce, y aumenta la tasa de ahorro. Por ello, para mantener el equilibrio de pleno empleo, es necesario que el porcentaje sobre PIB de los otros componentes de la demanda agregada aumente. No obstante, en tanto la inversión depende principalmente de los \textit{animal spirits}, lo preferible será aumentar el gasto público. Basándose en esto, algunos economistas predijeron que la economía experimentaría lo que llamaron \textbf{estancamiento secular} (Hansen, 1939), entendida como una larga depresión de duración indefinida, a menos que el sector público utilizara la política fiscal para incentivar la demanda agregada.
}

Si bien es cierto que esta concepción del destino de las economías avanzadas cayó en el olvido rápidamente, SUMMERS (2014) resucitó estas conjeturas y las aplicó a la situación económica actual de los países avanzados (Japón, Estados Unidos y Europa), afirmando que: \textit{Puede ser imposible que la economía alcance simultáneamente el pleno empleo, un crecimiento satisfactorio y la  estabilidad financiera simplemente mediante la operación de la política monetaria convencional}.\footnote{%
    SUMMERS justificaba su observación en que durante los años de bonanza desde el pinchazo de la burbuja Puntocom a principios de los años 2000's y el comienzo de la crisis subprime en 2007, la economía de Estados Unidos, aún teniendo unos tipos de interés muy bajos y generando una burbuja en los mercados de activos de enormes proporciones, no experimentó presiones inflacionistas, no logró bajar la tasa de desempleo sustancialmente por debajo de la media histórica y tampoco logró crecimientos significativamente por encima de la tendencia.
}

Para hacer frente a estos impactos negativos del estancamiento secular, SUMMERS plantea dos estrategias:\footnote{%
    En su artículo \textit{U.S. Economic Prospects: Secular Stagnation, Hysteresis, and the Zero Lower Bound} SUMMERS considera una tercera estrategia basada en la paciencia: en ocasiones las políticas económicas tienen limitado impacto y esto puede revertirse en el futuro. El autor sugiere que esta parece ser la estrategia seguida por países como Japón.
} 
\begin{lnum}
    \item \textbf{Reducción del Tipo deinterés real}. Buscar alternativas para reducir aún más las tasas de interés reales. Por ejemplo, estableciendo un objetivo de tasa de inflación más alto. No obstante, estas estrategias conllevan el riesgo de que incluso si aumentan el nivel de demanda, también es probable que aumenten los riesgos para la estabilidad financiera, lo que a su vez puede tener consecuencias negativas sobre la demanda; y, 
    \item \textbf{Aumento de la inversión}. Aumentar la demanda mediante el incremento de la inversión y la reducción del ahorro. Por ejemplo, aumentando la inversión pública, reduciendo las barreras sobre la inversión privada o fomentando las exportaciones.
\end{lnum}

El contexto económico asociado a un régimen de estancamiento secular tiene implicaciones importantes en términos de Política Económica. En relación a la política monetaria, su eficacia para estimular el consumo y la inversión es muy reducida cuando los tipos de interés iniciales están cerca del ELB. Este déficit de demanda nacional se manifestaría en movimientos internacionales de capital que podrían acentuar los desequilibrios exteriores. En principio, cabría esperar que los países abocados a un mayor envejecimiento exportarán capital, de manera que aumentarían los superávits de sus balanzas por cuenta corriente, mientras que en los países con poblaciones más jóvenes tendrían balanzas por cuenta corriente deficitarias. Esta salida de capitales hacia los mercados emergentes llevaría asociada una disminución de los tipos de cambio reales para los países industriales, una mayor competitividad y una mayor demanda de exportaciones. Adicionalmente, cabe esperar migraciones internacionales desde los países con poblaciones más jóvenes hacia aquellos otros con poblaciones más envejecidas, si bien todos los países experimentarán una escasez relativa de trabajadores jóvenes.

\subsection{Perturbaciones de Oferta}
Por otro lado, como ya anticipámos al principio, encontramos las perturbaciones cuyo origen se encuentra en la oferta. Entre los más comunes, cambios exógenos en la productividad o tecnología de las empresas y en los costes de producción.\footnote{%
    Las crisis que siguieron a los fuertes incrementos del precio del petróleo experimentados en 1973 y 1979, pusieron de manifiesto la incapacidad de la Política Económica para dar respuesta a este tipo de perturbaciones. Hasta entonces, esta se basaba en el manejo de una Curva de PHILLIPS tradicional y unos modelos que, entre otros aspectos, no incorporaban cambios en las expectativas de los agentes. Es en este marco cuando la teoría de los ciclos económicos reales pasa a ser el paradigma bajo el cual se concibe la política económica de la época, enfatizando el papel de las perturbaciones tecnológicas y factores de oferta en la generación de ciclos económicos. Su concepción sobre la intervención se consideraba como un agravante de las fluctuaciones, ya que se entendía que éstas respondían a ajustes en los mercados necesarios para mantener el equilbrio entre oferta y demanda. 

Como se ha dicho, este tipo de perturbaciones dejaron de ser particularmente immportantes a partr de los años 80's, y ante la incapacidad del modelo TCR de explicar los eventos económicos que se sucedían, la teoría NEK-2 se estableció, y ha dominado el pensamiento de política económica desde los años 90's.
}

Si bien las perturbaciones de oferta no han sido las predominantes en la generación de ciclos económicos desde la segunda mitad del siglo XX, es importante conocer su funcionamiento, pues cuando impactan, las consecuencias económicas y mecanismos de transmisión difieren de los de shocks de demanda y por lo tanto, requieren una respuesta distinta por parte de la Política Económica.

\subsubsection{Shock (positivo): Productividad Total de los Factores (tecnología)} 
Comenzaremos por analizar cambios en la productividad total de los factores.\footnote{%
    Recordemos que la PTF se suele inedir como el residuo de Solow, es decir, todo aquello que afecta al producto y que no se recoge en los factores productivos capital y trabajo. Si bien existen numerosos refinamientos del residuo de Solow (por ejemplo. incluyendo el capital humano como explicativa). esta variable no deja de ser un cajón de sastre que recoge multitud de factores que acaban afectando al producto final.
} 

Supongamos una perturbación positiva de oferta que se traduce en un aumento de la productividad total de los factores ($a_t$) de manera permanente y no anticipada. ¿Cuál es el mecanismo de transmisión al resto de la economía? Existen diversos efectos derivados de dicho shock pero lo que cabe destacar es que, en el modelo presentado bajo la regla monetaria de inflation targeting: \textbf{no se genera ineficiencia, por lo que no hay desplazamiento alguno}. No obstante, podemos resumir los efectos dinámicos de la siguiente forma:\footnote{%
    Directamente, $a_t$ entra en cuatro sitios: en la producción agregada $y_t=a_t+n_t$; en el coste marginal con coeficiente $-(1+\varphi)$; en el output natural con coeficiente $\Gamma=(1+\varphi)/(\sigma_\alpha+\varphi)$; y en el tipo natural a través del término proporcional a $-\sigma_\alpha\Gamma(1-\rho_a)$. Indirectamente, por moverse $mc_t$, $\overline{y}_t$ y $\overline{rr}_t$, se pueden ver alterados $\pi_{H,t}$, $x_t$, el tipo de interés nominal bajo una regla monetaria, el tipo de cambio real, los términos de intercambio $s_t$ y las exportaciones netas $nx_t$. 

La fuerza de cada canal depende de $\varphi, \sigma, \psi, \omega$ y $\rho_a$. Un mayor $\varphi$ hace que el shock reduzca más el coste marginal y eleve más el output natural. Un mayor $\psi$ reduce $\sigma_a$, amortiguando el traspaso de productividad a output natural. Un mayor $\omega$ amplifica el canal de términos de intercambio hacia demanda externa. Una mayor persistencia $\rho_a$ cambia la respuesta del tipo natural y, con ello, la parte intertemporal de la transmisión.
}
\begin{lnum}
    \item \textbf{Output}. El primer canal de transmisión del que cabe hablar es el de la producción: un aumento de $a_t$ eleva directamente la producción por unidad de trabajo ($y_t=a_t+n_t$). Además, manteniendo constante el resto, un aumento de $a_t$ reduce el coste marginal con elasticidad directa $\partial mc_t/\partial a_t=-(1+\varphi)$. Esta es la primera relación estructural importante: el shock de PTF es desinflacionario por el lado de costes;
    \item \textbf{Output potencial}. Además, un shock positivo de PTF eleva el output natural: mayor cuanto mayor sea $1+\varphi$ y menor cuanto mayor sea $\sigma_\alpha+\varphi$, $\Gamma=\frac{1+\varphi}{\sigma_\alpha+\varphi}$. Esta fracción es fundamental para política económica porque dice cuánto del aumento de productividad se traduce en capacidad productiva sin tensiones inflacionarias. En una economía más abierta, $\alpha$ es mayor y, por tanto, $\sigma_\alpha+\varphi$ también. Luego, ceteris paribus, una mayor apertura amortigua el impacto de $a_t$ sobre el output natural. La intuición es que, \textbf{en economía abierta, una parte mayor del ajuste pasa por los términos de intercambio y por la demanda relativa, no solo por la expansión de producción doméstica};
    \item \textbf{Inflación: CP}. Como $a_t$ empuja al alza $\bar y_t$, si el output observado $y_t$ no aumenta en la misma cuantía de forma inmediata, la brecha $x_t$ cae. Por tanto, incluso aunque el shock sea expansivo en términos de output potencial, puede ser contractivo para la brecha de producción y, en consecuencia, desinflacionario: un shock de oferta positivo suele elevar el primero y reducir el segundo;
    \item \textbf{Output gap}. Sin embargo, esto exige hablar del comportamiento del output gap. $a_t$ afecta indirectamente a la demanda porque mueve el tipo natural de interés: $\partial \overline{rr}_t/\partial a_t=-\sigma_\alpha\Gamma(1-\rho_a)$. Esto implica que, para un shock tecnológico transitorio $a_t$ con persistencia $\rho_a<1$, el tipo natural se mueve con la productividad: la PTF altera el perfil deseado de consumo y producción sin rigideces, y por esa vía desplaza $\overline{rr}_t$. Lo que nos conduce al elemento clave: (i) \textbf{si la autoridad monetaria acompañaa el tipo real de mercado con ese nuevo natural, la transición no genera ineficiencias}; (ii) si la Política Monetaria no acomoda, aparece un gap de demanda ($x_t$) y una desinflación;
    \item \textbf{Demanda externa}. Por último, la caída del tipo real relativo influye en el tipo de cambio y en los términos de intercambio. En una pequeña economía abierta, ese canal externo importa porque los términos de intercambio se mueven con los diferenciales de tipos reales esperados. Dada la constancia del tipo de interés nominal mundial, la paridad descubierta de intereses implica una depreciación nominal inicial seguida de expectativas de apreciación futura, tal como refleja la respuesta del tipo de cambio nominal. En comparación con otros regímenes, la constancia de los precios domésticos da lugar a una \textbf{mayor depreciación real} y, por tanto, a una expansión adicional de la demanda y del output a través del \textbf{aumento de las exportaciones netas}. Dado que los precios mundiales son constantes y los términos de intercambio son estacionarios, la constancia de los precios domésticos implica una respuesta del tipo de cambio nominal con reversión a la media.
\end{lnum}

En términos gráficos, lo que vemos es que, ante reacciones adecuadas y rápidas de la Autoridad Monetaria, las dinámicas serán las siguientes:

\imagenfit{ShockProductividad_Positivo.png}{Respuestas dinámicas ante un shock ante Políticas Monetarias alternativas}{GALÍ y MONACELLI (2005)}

También observamos que el shock provoca una reducción persistente del tipo de interés doméstico, ya que es necesaria para sostener la expansión transitoria del consumo y del output consistente con la asignación de equilibrio bajo precios flexibles. Obsérvese que en los resultados expuestos se comparan las dinámicas teniendo en cuenta diferentes regímenes.\footnote{%
    Del artículo de los autores:

\textit{En primer lugar, que ambas reglas generan, a diferencia de la política óptima, una caída permanente tanto de los precios domésticos como del IPC. La raíz unitaria en los precios domésticos se refleja entonces, bajo ambas reglas, en una raíz unitaria en el tipo de cambio nominal.

Una diferencia clave entre las dos reglas de Taylor se refiere al comportamiento de los términos de intercambio. Mientras que bajo DITR los términos de intercambio se deprecian inicialmente y comienzan inmediatamente a revertir hacia su estado estacionario (de forma muy similar a la política óptima), bajo CITR la respuesta inicial es más atenuada y va seguida de una dinámica en forma de joroba. La intuición es sencilla. Bajo ambas reglas, el aumento de la productividad doméstica y la depreciación real necesaria conducen, para unos precios domésticos dados, a un incremento de la inflación del IPC. Sin embargo, bajo CITR la estabilización deseada de la inflación del IPC se logra en parte, en comparación con DITR, mediante una respuesta más moderada de los términos de intercambio (ya que estos afectan al IPC) y una caída de los precios domésticos. Esta caída de precios requiere, a su vez, una brecha de producción negativa y, por tanto, una política monetaria más contractiva (es decir, un tipo de interés más elevado). Bajo nuestra calibración, esto se traduce en un aumento inicial tanto del tipo de interés nominal como del real, manteniéndose posteriormente el tipo real sistemáticamente por encima del implicado por la política óptima o por una regla DITR.

Por último, la misma figura muestra las funciones impulso-respuesta correspondientes bajo un régimen de tipo de cambio fijo (PEG). Obsérvese que las respuestas de la brecha de producción y de la inflación son cualitativamente similares a las del caso CITR. Sin embargo, la imposibilidad de reducir el tipo de interés nominal y permitir la depreciación de la moneda —como sería necesario para sostener la expansión del consumo y del output requerida para replicar la asignación de precios flexibles— conduce a una respuesta muy limitada de los términos de intercambio y, en consecuencia, a una amplificación de las respuestas de la inflación doméstica y de la brecha de producción. De manera interesante, bajo un PEG, la completa estabilización del tipo de cambio nominal genera estacionariedad en el nivel de precios domésticos y, a su vez, también en el nivel del IPC (dada la estacionariedad de los términos de intercambio). Esta es una propiedad que el régimen de tipo de cambio fijo comparte con la política óptima descrita anteriormente. La estacionariedad del nivel de precios explica también por qué, en respuesta al shock, la inflación doméstica cae inicialmente y posteriormente aumenta de forma persistente por encima de su nivel de estado estacionario.}
}

\paragraph{Implicaciones de Política Económica}
La principal implicación de Política Económica se puede ver fácilmente con la figura anterior: un shock positivo de PTF es, de entrada, desinflacionario y eficiente, porque mejora la frontera productiva y, con una reacción adecuada bajo la política óptima (objetivo estricto de inflación) el banco central puede replicar la asignación flexible estabilizando $\pi_{H,t}=0$; en ese caso también se estabiliza $x_t=0$.\footnote{Cuando se logra neutralizar la distorsión monopolística con subsidio al empleo.
}

Por tanto, la inflación doméstica y la brecha de producción permanecen en cero, y el output de la pequeña economía abierta aumenta siempre en respuesta a un shock tecnológico positivo doméstico. Es decir, la política monetaria óptima no aplana el efecto real beneficioso de la productividad; lo que \textbf{elimina es la distorsión nominal}.  


\subsubsection{Shocks (negativos): Costes de producción}
Este modelo también permite el análisis de otro tipo de perturbaciones que surgen por el lado de la oferta y que afectan directamente a la inflación: los costes de los insumos productivos.

Si bien es cierto que el modelo original (así como la extensión propuesta por FUNARTELLO) no contienen este tipo de shocks, podemos introducirlo fácilmente reformulando la Curva de PHILLIPS introduciendo un término de error:

\eqblock{
\begin{aligned}
    &\text{Curva de PHILLIPS}:&&\pi_{H,t}=\beta \mathbb E_t[\pi_{H,t+1}]+\kappa_\alpha x_t+u_t,\quad\text{donde}\quad u_t\sim N(0,\sigma_u)\\[0.5em]
    &\kappa_\alpha\to0\Longrightarrow &&|Var(\pi_{H,t})|\gg0\iff |Var(u_t)|>0\\
    &\Longrightarrow \kappa_\alpha\equiv \lambda(\sigma_\alpha+\varphi):&&\sigma_\alpha=\frac{\sigma}{1-\alpha+\alpha\omega}\Longrightarrow\;(\uparrow\alpha\Rightarrow\;\downarrow\kappa_\alpha)
\end{aligned}
}{Curva de PHILLIPS: Inflación de costes. Perturbaciones: shocks (negativo) de costes}

Este nuevo término de error recoge una perturbación exógena del margen o del coste que altera la relación entre inflación y output gap para un mismo nivel de equilibrio (\textit{ceteris paribus}). Por eso, como veremos, es especialmente importantes para la Política Monetaria.\footnote{%
    Sutherland lo formula precisamente así: son perturbaciones, como cambios en markups o en impuestos distorsionantes, que mueven precios sin implicar un cambio en el nivel socialmente óptimo de output.

Recordemos que en el modelo básico sin cost-push, cuando neutralizamos la distorsión monopolística con el subsidio al empleo ($\tau$), estabilizar la inflación doméstica equivale a estabilizar el output gap. Pero con $u_t\neq 0$, eso deja de ser cierto. Si el banco central mantiene un objetivo estricto de inflación entonces de la NKPC se sigue: $x_t=-\frac{u_t}{\kappa_\alpha}$.

Luego, ante un shock negativo en costes $u_t<0$, el output gap debe ser positivo para sostener el objetivo de inflación cero. Es decir, el banco central necesita acomodar el shock favorable con una política más expansiva, de forma que el mayor empuje desinflacionario del lado de costes quede compensado por una mayor demanda. Si no acomoda, la inflación cae por debajo del objetivo. 

En esta nueva formulación, la variable determinante será $\kappa_\alpha$: cuanto menor sea $\kappa_\alpha$, mayor tendrá que ser el movimiento del output gap necesario para neutralizar un mismo shock de costes. En otras palabras, \textbf{el coste real de estabilizar inflación frente a un shock de costes es inversamente proporcional a la pendiente} de la curva de PHILLIPS.
} 

Si suponemos un shock positivo ($u_t>0$) que desplace la NKPC hacia arriba, el Banco Central aumentará el tipo nominal para generar un aumento del tipo real relativo al natural, ante las expectativas de inflación. Ese movimiento provocará en la IS abierta una contracción de la demanda agregada que no se agota en consumo interno: el aumento del tipo de interés afecta al tipo de cambio nominal esperado vía paridad no cubierta de intereses (PNCI), mueve el tipo de cambio real y los términos de intercambio, y eso repercute sobre exportaciones netas. Por tanto, el shock de costes se transmite a la demanda a través de la reacción monetaria y del canal externo. Las variables directamente afectadas son, por tanto:
\begin{lnum}
    \item \textbf{Inflación}. Un aumento de la inflación, más pronunciado cuánto menor sea el parámetro de apertura ($\kappa_\alpha$);
    \item \textbf{Output y ouput-gap}. Una contracción de la demanda agregada vía canal nominal por la reacción del Banco Central ante un aumento de la inflación;
    \item \textbf{Términos de intercambio y BC}. Y, por último, relativo al canal exterior la reacción de la Autoridad Monetaria afectará doblemente: (i) por un lado vía términos de intercambio (por PNCI), el tipo de cambio nominal y el tipo de cambio efectivo real se apreciarán; (ii) por lo que respecta a la balanza comercial, deteriorando su saldo. 
\end{lnum}

Las variables paramétricas que gobiernan la intensidad del mecanismo son $\kappa_\alpha, \sigma_\alpha$, el grado de apertura $\alpha$, y la elasticidad de sustitución entre bienes nacionales y extranjeros ($\eta$). Cuanto más abierta es la economía, más pequeño es el valor de $\kappa_\alpha$, más plana es por tanto la Curva de PHILLIPS y, ante un shock de costes: (i) más importante es el canal de precios relativos; y, (ii) más delicado (complejo) se vuelve el trade off entre estabilización de inflación y de actividad. 

Además, démonos cuenta de que lo que cambia respecto a un shock de demanda es que aquí la perturbación original entra por la oferta nominal y solo después pasa a la demanda mediante la regla monetaria.

\paragraph{Implicaciones de Política Económica}
No obstante, la implicación más importante de Política Económica es precisamente la inexistencia de la divina coincidencia: si fijamos $\pi_{H,t}=0$, aceptamos movimientos del output gap; si intentamos suavizar el output gap, aceptaremos desviaciones de inflación. Por tanto, haga lo que haga la autoridad monetaria, cualquier shock de costes que experimente la economía no podrá ser acomodado con su actuación y deberán buscarse alternativas. Situación que pone de manifiesto, al igual que lo visto en los shocks de preferencias, tanto los límites de la Política Monetaria como los límites del modelo.

Por otro lado, algunos autores argumentan que la existencia de push costs hace que la política óptima deja de ser de objetivo inflacionario y pasa a parecerse a un \textit{flexible inflation targeting}.


\clearpage
\section{Nuevas lineas teóricas: MABM y Hechos Estilizados}
A lo largo de nuestra exposición hemos visto la modelización tradicional de la oferta y demanda agregada a través de los modelos NOEM, que son el caballo de batalla de la mayoría de economistas en la actualidad. 

Además, hemos visto cómo estos proporcionaban indicaciones de Política Económica como respuesta a shocks, principalmente centrados en la actuación de la Política Monetaria (como ha sido tradición desde los años 70's) y relegando a un lugar muy secundario, casi irracionalizando las respuestas de Política Fiscal. Todo esto no es más que un resultado directo de la aproximación metodológica que se ha seguido. 

Cualquier modelo de optimización dinámica puede hacerse \textit{ad hoc} para mostrar, con mayor o menor rigor, lo que el investigador pretende decir. No obstante, esto no es una limitación objetivo, sino más bien subjetivo de éstos.

Ahora bien, para encontrar una limitación más que suficiente para buscar otras aproximaciones metodológicas solo es necesario recordar cómo son tratadas las perturbarciones que azotan constantantemente nuestras economías lejos del \textit{equilibrio}. La propia etimología de la palabra \textit{shock} parece justificar que las economías son \textit{empujadas}, \textit{golpeadas} o \textit{azotadas} fuera del equilibrio por un agente totalmente externo. Más aún, el equilibrio global sufre esta misma dolencia. Pero, ¿tiene esto algún sentido? ¿quiere Dios atormentarnos por nuestros pecados?

La verdad es, que por mucho ajuste empírico que acabe teniendo (sin duda muchas veces debido a la gran cantidad de cambios \textit{ad hoc} que se realizan sobre estos), la forma en la que se trata de la dinámica económica dista mucho de lo que cualquier persona consideraría \textit{razonableme}. 

\textbf{¿Qué queremos decir al final del apartado?} Por veremos brevemente otra linea metodológica alternativa, la de los MABM, que pretende representar dinámicamente la economía mediante simulaciones y no mediante optimización. Veremos como esto nos permite un ajuste ex ante mucho más correcta, y parece reproducir características/comportamientos ex post singulares y muy relevantes, puesto que son estos los que precisamente atormentan con más frecuencia la economía global y el conjunto de sus dirigentes. Además, veremos como la evidencia empírica ha hecho una gran labor de recaptación, formando un conjunto de hechos estilizados especialmente relevantes. 

\subsection{MABM: Mercados Financieros}
Para presentar un ejemplo claro de metodología por simulación recurrimo a un Modelo Macroeconómico Basado en Agentes que pretende simular las fluctuaciones en los mercados financieros internacionales. En concreto, en el mercado FOREX, sin duda el más importante en volumen de todos.\footnote{%
    [Ver Tema 3.B.14]. Todos los modelos de fijación del Tipo de Cambio dentro del marco metodológico clásico construyen este como un precio compuesto de diferentes variables. Por tanto, subyace la idea de que el Tipo de Cambio tiene un precio/valor de equilibrio o natural y que las desviaciones de este valor pueden atribuirse a shocks estocásticos, cambio de los determinantes o, en última instancia, mala especificación o base teórica. (con la única excepción, aunque tampoco del todo , del Modelo de Agentes Heterogeneos)

Ahora bien, ¿y si el Tipo de Cambio no fuese más que el resultado de una serie de comportamientos endógenos que poco, o nada, tienen que ver con éste pero en última instancia emerge como una propiedad del sistema? Es decir, ¿y si no existiera un equilibrio sino, más bien, un resultado en constante mutación/variación/alteración? Este es precisamente el enfoque que toman los Modelos Basados en Agentes, que dicen ser \textit{bottom-up}.

El más célebre de estos en la materia es el recientemente publicado por un grupo interuniversitario asociado a la Universidad de Pise (\href{https://papers.ssrn.com/sol3/papers.cfm?abstract_id=4903314}{\textit{The interplay between real and exchange rate market: an agent-based model approach}. DOMENICO et al., 2024}), referente en este campo.
}

\subsubsection{Tipo de Cambio: Naturaleza dual}
Supongamos que el Tipo de Cambio tiene una naturaleza dual: refleja fundamentos reales pero también representa un activo financiero determinado por decisiones especulativas de agentes con expectativas heterogéneas (fundamentalistas y chartistas). 

\subsubsection{Desarrollo formal: Multi-país, multi-sector y multi-agente}
DOMENICO et al., (2024)\footnote{%
    Multi-país y multi-sector, que integra simultáneamente comercio internacional, innovación, expectativas heterogéneas en el mercado de divisas y la actuación del Banco Central, extendiendo el marco de DOSI et al. (2019).
} construyen el siguiente marco simulado:

\imagenfit{MABM_2024_EsquemaEcoReal.png}{Proceso de interacción entre la economía real y el comportamiento financiero del tipo de cambio. Dinámica multi-país, multi-sector y multi-agente}{\href{https://papers.ssrn.com/sol3/papers.cfm?abstract_id=4903314}{\textit{The interplay between real and exchange rate market: an agent-based model approach}. DOMENICO et al. (2024)}}

Marco en el cual, los agentes que lo conforman deben tomar una serie de decisiones en cada periodo y se suceden una serie de eventos:
\begin{lnum}
    \item Las empresas de las industrias de bienes de consumo realizan I+D para descubrir nuevas técnicas e imitar a los competidores más cercanos a la frontera tecnológica (análisis envolvente de datos: DEA). Las empresas pueden mejorar su productividad laboral si la innovación o la imitación tienen éxito;
    \item Tienen lugar las decisiones de producción y empleo. Las empresas de bienes de consumo fijan su producción deseada dado su nivel de demanda esperada. En consecuencia, contratan trabajadores y amplían su capacidad productiva si es necesario;
    \item Se fijan los salarios monetarios y los precios;
    \item El tipo de cambio bilateral viene determinado por las decisiones de comportamiento de los operadores comerciales en el mercado de divisas;
    \item Se abren los mercados internacionales de bienes de consumo con competencia imperfecta. Los trabajadores gastan sus ingresos en bienes nacionales e importados. Las cuotas de mercado de las empresas evolucionan según su competitividad en precios;
    \item Tiene lugar la entrada y salida de empresas. Nuevas empresas sustituyen a las empresas con cuotas de mercado casi nulas que abandonan el mercado.
\end{lnum}

Al final de cada paso de tiempo, se calculan las variables agregadas (por ejemplo, producción, consumo, exportaciones, importaciones, etc.), sumando las correspondientes variables microeconómicas.\footnote{%
    Para ver las normas/procesos de decisión; las probabilidades asociadas y las dinámicas de cada etapa consultar el paper.
}

\imagenfit{Dinamica_MABMEcoFinanciera.png}{Diagrama de flujo que ilustra los principales pasos diferenciadores entre el modelo de referencia (blanco + azul) y el modelo con comportamiento especulativo (amarillo + azul) y el Banco Central (verde)}{\href{https://papers.ssrn.com/sol3/papers.cfm?abstract_id=4903314}{\textit{The interplay between real and exchange rate market: an agent-based model approach}. DOMENICO et al. (2024)}}

Una particularidad (notable) del modelo es que no permite soluciones \textit{closed-form}, común de los modelos basados en agentes.\footnote{%
    Esta limitación surge de las no linealidades inherentes en las reglas de decisión de los agentes y en los patrones de sus interacciones. En consecuencia, para investigar la dinámica de las variables micro y macro, es necesario recurrir a simulaciones. Se imponen parámetros estructurales idénticos en torno a países, industrias y empresas. La dinámica evolutiva es un resultado endógeno.
}

\subsubsection{Resultados: Implicaciones de Política Económica}
Las simulaciones muestran, como resultado principal, que la introducción de comportamiento heterogeneidad de agentes, con diferentes objetivos, permite reproducir de manera conjunta una amplia batería de hechos estilizados empíricos en las fluctuaciones cambiarias: colas gruesas, clustering de volatilidad, efectos de apalancamiento, presencia de raíz unitaria, burbujas y crashes en el tipo de cambio, así como el conocido \textit{disconnect puzzle} entre tipos de cambio y fundamentos. 

En términos de Política Económica, las conclusiones son también infinitamente más esclarecedoras puesto que el marco de simluación puede alterarse adecuandolo a las posibilidades reales del mercado FOREX. Es decir, no se realizan ajustes \textit{ad hoc} sino que se plantean reformas estructurales que permiten ser simuladas y evaluar este impacto en consecuencia. 

Por ejemplo, en este caso particular demuestran como los chartistas (seguidores de tendencias) son un factor clave de inestabilidad ya que aumentan la duración y la amplitud de los ciclos financieros y, a través del canal cambiario, amplifican las fluctuaciones reales, afectando al crecimiento, al comercio y a la sincronización internacional de los ciclos.

Por lo que respecta la actuación de Políctica Monetaria, los autores demuestran como la participación activa del Banco Central en el mercado de divisas, puede reducir la amplitud de los ciclos y limitar la especulación mediante intervenciones esterilizadas del tipo \textit{leaning against the wind}. Sin embargo, también subrayan con fuerza que su eficacia es limitada y dependiente del diseño: las intervenciones continuas logran frenar burbujas, pero pueden aumentar la volatilidad real; las intervenciones esporádicas con objetivos amplios resultan prácticamente ineficaces. El banco central, por tanto, enfrenta trade offs inevitables y no puede eliminar la inestabilidad endógena generada por expectativas heterogéneas.

En conclusión, el paper sostiene que la inestabilidad financiera y cambiaria no es exógena ni accidental como sería felizmente tratada en el marco expuesto, sino una \textbf{propiedad emergente} de economías abiertas con agentes heterogéneos y comportamiento especulativo. 

\subsection{Evidencia Empírica: Algunas conclusiones}
Hasta ahora hemos visto que la principal respuesta de los países ante perturbaciones es a través del canal de tipos de interés. Este tipo de respuesta no se queda restringida a las fronteras de un país, sino que en el caso de economías abiertas, se transmiten entre países. Muchos han sido los trabajos que han analizado la transmisión internacional de shocks entre países, entre ellos de cambios en los tipos de interés. Sin embargo, tras la crisis financiera de 2008 se incluyó en este mecanismo de transmisión los mercados financieros y flujos de capital como potenciadores de dicho mecanismo. Es lo que se denomina el Ciclo Financiero Global (Rey, 2015). Además, tras la crisis financiera la transmisión de estas perturbaciones fue incluso más importante, no sólo entre países avanzados, sino a los países emergentes. 

\subsubsection{Ciclo Financiero Global}
El Ciclo financiero global se caracteriza por fuertes movimientos conjuntos de la actividad a través de diversas variables a nivel internacional. 

En épocas de bonanza, un nivel de actividad económica alto genera inversiones y empleo. Muchas de estas inversiones se financian con crédito, abundante durante las expansiones. Esto puede llevar a la generación de situaciones de riesgo (préstamos concedidos que en otras condiciones serían demasiado arriesgados). 

Este mayor riesgo asumido por los inversores puede generar inestabilidad financiera y macroeconómica en el momento en que la expansión pierda fuerza. De hecho, esto conduce a REY (2015) a cuestionar el trilema de la Política Monetaria, argumentando que en realidad se debería hablar de dilema ya que la clave radica en hasta qué punto se puede mantener la independencia en política monetaria con libre movilidad de capitales. 

La existencia de un ciclo financiero global hace que decisiones de Política Monetaria en un país generen cambios similares en otros países, debido precisamente a la libre movilidad de capitales. En este escenario las decisiones tomadas por la Reserva Federal de Estados Unidos son clave en la generación de estos ciclos. Varios son los canales por los que la política monetaria de la Fed puede tener implicaciones a nivel de ciclo financiero global:
\begin{la}
    \item \textbf{Acelerador financiero: Valoración de activos}. Por un lado, los efectos indirectos de variaciones de los tipos de interés en Estados Unidos vía la valoración de activos ligados al dólar estadounidense y al flujo de capitales a nivel internacional; y,
    \item \textbf{Riesgo}. Por otro lado, las posiciones de riesgo asumidas por los inversores se pueden ver afectadas por las decisiones de la Fed. Tipos de interés más o menos altos pueden generar flujos de capitales en búsqueda de mayores rendimientos. 
\end{la}

Además, la estabilidad financiera global también se puede ver afectada por decisiones tomadas por la Fed. La Gran Recesión puso de manifiesto la fuerza del ciclo financiero global en la transmisión de perturbaciones a nivel internacional. Desde entonces, la coordinación de las políticas monetaria y fiscal a nivel nacional así corno las monetarias a nivel internacional han sido continuas. Adicionalmente, la política macroprudencial ha surgido como un nuevo instrumento para las autoridades de las economías avanzadas con el fin de mantener la estabilidad financiera.\footnote{%
    Es importante destacar que el uso de la política macroprudencial estaba ampliamente extendido entre los países emergentes más vulnerables a cambios en las condiciones financieras internacionales.
}

\subsubsection{Economía real: Shocks de costes}
Otra oprtunidad de hora que se le proporcionó a los investigadores empíricos para tratar de estudiar las interrelaciones y efectos que tenía sobre una economía algún fenómeno exógeno fue precisamente las sucesivas crisis del petroleo durante los años 70's. 

Aunque no fue la última, si duró lo suficiente como para poder estudiar el impacto sobre la economía real de una perturbación de costes. Particularmente para los países más desarrollados, donde esta materia prima es de especial interés ya que dependen en su mayoría de ella. Incluso existe una amplia literatura (por ejemplo Hamilton, 1983 y 1996) que documenta cómo la mayoría de las recesiones sufridas en Estados Unidos vinieron precedidas de aumentos en el precio del petróleo. 

PEERSMAN y VAN ROBAYS (2009) analizan los efectos macroeconómicos de perturbaciones del precio del petróleo en Europa distinguiendo el origen de las mismas y los comparan con los efectos en Estados Unidos. Los autores destacan la importancia de distinguir entre aumentos de los precios del petróleo debidos a falta de oferta, a aumentos de la demanda de petróleo o por aumentos globales en el nivel de actividad económica. Además, distinguen entre efectos de primera y segunda ronda:
\begin{la}
    \item \textbf{Efectos primera ronda}. Los efectos de primera ronda comprenden el aumento de precios vía aumento de costes marginales y son los que dominan en la transmisión en Estados Unidos; y,
    \item \textbf{Efectos segunda ronda}. Los efectos de segunda ronda incluyen aumentos en los salarios que contribuyen a una mayor inflación. Según los autores, estos efectos de segunda ronda son los que aumentan en Europa.
\end{la} 







\clearpage
\section{Conclusión} 

\subsection{Extensiones y relación con otras partes del Temario}


\subsection{Opinión}



\subsubsection{Idea final (Salida o cierra)}


\clearpage
\pagenumbering{gobble}
\MyBibliography
\MyBibItem{Autor}{Título}{Año}{DOI -- LINK}




%ANEXOS
\clearpage
\anexo{\textbf{Preguntas Test}}
%\input{\TemaCode/questions.tex} 


\clearpage
\anexo{Derivación del modelo GALI y MONACELLI (2005) -- NOEM}\label{anx:nek2_gali}

Para conocer la derivación completa del modelo, véase: \href{https://crei.cat/wp-content/uploads/users/pages/roes8739.pdf}{\textit{Monetary Policy and Exchange Rate Volatility in a Small Open Economy}. GALI, J. \& MONACELLI, T. (2005)}

Obviaremos todos los elementos relacionados con la maximización del consumidor y productor representativo, pasando a analizar directamente los resultados de ambos problemas y sus consecuencias a la hora de construir el Modelo de la NEK-2 para una economía pequeña y abierta.

\textsc{Enfoque de demanda (nacional e internacional)}

En primer lugar, del problema de optimización del consumidor representativo obtenemos las siguientes dotaciones de consumo óptimas para cada categoría de productos:

$$\begin{aligned}
    &\underbrace{C_{H,t}(j)=\left(\frac{P_{H,t}(j)}{P_{H,t}}\right)^{-\varepsilon}C_{H,t}}_{\text{\scriptsize Demanda de consumo óptima del bien (doméstico) $j$}};\qquad \underbrace{C_{i,t}(j)=\left(\frac{P_{i,t}(j)}{P_{i,t}}\right)^{-\varepsilon}C_{i,t}}_{\text{\scriptsize Demanda de consumo óptima del bien (importado) $j$}};\qquad \underbrace{C_{i,t}=\left(\frac{P_{i,t}}{P_{F,t}}\right)^{-\gamma}C_{F,t}}_{\substack{\text{\scriptsize Demanda de consumo óptima de bienes procedentes de $i$}\\\text{\scriptsize $\gamma$ mide grado de sustituibilidad entre diferentes $j$ extranjeros}}}\\[2em]
    &\text{donde,}\quad 
    \left\{\begin{aligned}
        &P_{H,t}\equiv\text{Índice de precios domésticos}\\
        &P_{i,t}\equiv \text{Índice de precios del país $i$}\\
        &P_{F,t}\equiv\text{Índice de precios importados}
    \end{aligned}\right.\\[1em]
    &\Longrightarrow\underset{\text{\scriptsize Dotaciones óptimas de consumo de bienes nacionales y bienes importados}}{\boxed{C_{H,t}=(1-\alpha)\left(\frac{P_{H,t}}{P_t}\right)^{-\eta}C_t\qquad \text{y,}\qquad C_{F,t}=\alpha\left(\frac{P_{H,t}}{P_t}\right)^{-\eta}C_t}}\\[1em]
    &\text{donde,}\quad \text{CPI}\;\;\equiv\;\;P_t=[(1-\alpha)(P_{H,t})^{1-\eta}+\alpha(P_{F,t})^{1-\eta}]^{\frac{1}{1-\eta}}\Longrightarrow
    \left\{\begin{aligned}
        &\alpha\equiv\;\;\substack{\text{\scriptsize Parámetro de sesgo nacional en preferencias}\\\text{\scriptsize p.t. proxy para grado de apertura}}\\[0.5em]
        &\eta\equiv\;\;\substack{\text{\scriptsize Mide el grado de sustituibilidad}\\\text{\scriptsize entre domésticos y extranjeros}}
    \end{aligned}\right.
\end{aligned}$$

Ahora que tenemos las demandas óptimas de consumo para cada nación, cada producto y de manera agregada en función del grado de sustituibilidad y del grado de apertura, podemos representar el CPI (log-linearizado) a través de los términos reale de intercambio. Por tanto, algunas variables especialmente interesantes para el posterior análisis se extraerán de las siguientes relaciones:

$$\begin{aligned}
    &\text{Términos de intercambio}:&&S_t=\frac{P_{F,t}}{P_{H,t}}\overset{\ln(\cdot)}{\Longrightarrow}s_t=\int_0^1s_{i,t}\mathrm di\Longrightarrow \underset{\substack{\text{\scriptsize Términos efectivos de intercambio (log)}\\\text{\scriptsize i.e. precios extranjeros en términos domésticos}}}{\boxed{s_t=p_{F,t}-p_{H,t}}}\\[1em]
    &\text{CPI}:&&p_t\equiv(1-\alpha)p_{H,t}+\alpha p_{F,t}\Longrightarrow {\boxed{p_t=p_{H,t}+\alpha s_t}}\\[1em]
    &\text{Inflación}: &&\pi_t=\underset{\equiv p_{H,t}-p_{H,t-1}}{\pi_{H,t}}+\alpha\Delta s_t\\[1em]
    &\overset{\substack{\text{\scriptsize Será crucial para determinar la Paridad No Cubierta de Intereses:}\\\text{\scriptsize Flujos de capital en estado estacionario}}}{\underset{\text{\scriptsize Si asumimos que se cumple la Ley del Precio Único: PPP}}{\text{Tipo de cambio efectivo nominal}}}&& e_t=\int_0^1e_{i,t}\mathrm di\Longrightarrow p_{F,t}=e_t+\underset{\substack{\text{\scriptsize World}\\\text{\scriptsize Price}\\\text{\scriptsize Index}}}{p_t^*}\Longrightarrow \boxed{s_t=e_t+p_t^*-p_{H,t}}\\[1em]
    &{\text{Tipo de intercambio real efectivo}:}&&q_t=(1-\alpha)s_t\iff \eta\neq 1
\end{aligned}$$

Por último, y como anticipábamos, bajo el supuesto de mercados financieros internacionales completos, log-linearizando las condiciones de los flujos de capital en torno a un estado estacionario de previsión perfecta y agregando sobre $i$, se obtiene la expresión habitual: 

$$r_t-r_t^*=\mathbb E_t[\Delta e_{t+1}]$$

Por tanto, combinando la definición de los términos de intercambio (en logaritmos), se obtiene la siguiente ecuación en diferencias estocástica: 

$$s_t\equiv (r_t^*-\mathbb E_t[\pi_{t+1}^*])-(r_t-\mathbb E_t[\pi_{H,t+1}])+\mathbb E_t[s_{t+1}]$$

Por tanto, los términos de intercambio quedan determinados de forma única en el estado estacionario de previsión perfecta. 

Este hecho, junto con el supuesto de estacionariedad de las perturbaciones del modelo y una normalización conveniente (que impone que la PPA se cumple en el estado estacionario), implica que:

$$\lim_{T\to\infty} E_t\{s_T\}=0$$ 

No obstante, aunque estas ecuaciones ofrecen una representación conveniente (e intuitiva) de la relación entre términos de intercambio y diferenciales de tipos de interés, no constituyen una condición de equilibrio independiente adicional.

\textsc{Enfoque de oferta (nacional)}

En el modelo, una empresa típica en la economía doméstica produce un bien diferenciado con una tecnología lineal representada por la función de producción: 

$Y_t(j)=A_t N_t(j)$, donde $a_t \equiv \log A_t\sim \text{AR (1)}$, $a_t=\rho_a a_{t-1}+\varepsilon_t, y j \in [0,1]$ es un índice específico de empresa. 

En consecuencia, el coste marginal real (expresado en términos de precios domésticos) será común a todas las empresas nacionales y viene dado por:

$mc_t = -\nu + w_t - p_{H,t} - a_t$, donde $\nu \equiv -\log(1-\tau)$, siendo $\tau$ un subsidio al empleo.

Si consideramos $Y_t$ un índice de producción agregada doméstica, análogo al introducido para el consumo. Y, $N_t \equiv \int_0^1 N_t(j) dj = \frac{Y_t Z_t}{A_t}$.\footnote{Donde $Z_t \equiv \int_0^1 \frac{Y_t(j)}{Y_t} \mathrm dj$.
}

Por tanto, hasta una aproximación de primer orden, se obtiene la \textbf{relación agregada}: 

$$y_t = a_t + n_t$$

Por otro lado, considerando fijación de precios à la CALVO: 

$$\bar p_{H,t} = \mu + (1-\beta\theta)\sum_{k=0}^{\infty}(\beta\theta)^k\mathbb E_t\{mc_{t+k} + p_{H,t}\}$$

Donde $\bar p_{H,t}$ denota el logaritmo del nuevo precio fijado y $\mu \equiv \log\left(\frac{\varepsilon}{\varepsilon-1}\right)$, que corresponde al logaritmo del mark-up en el estado estacionario.\footnote{%
    Así, vemos que la decisión de fijación de precios en este modelo (al igual que en su equivalente de economía cerrada) es de carácter prospectivo (forward-looking). La razón es sencilla: las empresas que ajustan precios en un periodo dado saben que el precio fijado permanecerá vigente durante un número (aleatorio) de periodos. Como resultado, fijan el precio como un margen sobre una media ponderada de los costes marginales futuros esperados, en lugar de basarse únicamente en el coste marginal corriente. Obsérvese que, en el límite de precios flexibles (es decir, cuando $\theta \to 0$), se recupera la regla habitual de margen: $\bar p_{H,t} = \mu + mc_t + p_{H,t}$.
}

\textsc{Equilibrio: Ecuaciones del Modelo NEK-2 (NOEM)}

Por último, e introduciendo las condiciones de equilibrio en nuestro modelo, log-linearizando en estado estacionario, obetenemos las siguentes ecuación relativa a la producción-consumo de equilibrio:

$$\begin{aligned}
    &\text{Curva IS}:&&y_t=c_t+\alpha\gamma s_t+\alpha\left(\eta-\frac{1}{\sigma}\right)q_t\\
    &&&\substack{\text{si, }\quad \omega\equiv\sigma\gamma+(1+\alpha)(\sigma\eta-1)}{\Longrightarrow}\boxed{y_t=c_t+\frac{\alpha\omega}{\sigma}s_t\quad (1)}\overset{\text{\scriptsize agregando}}{\Longrightarrow}\boxed{y_t^*=c_t^*\quad (2)}\\[1em]
    &&&\text{Combinando (1) y (2)}\Longrightarrow y_t=y_t^*+\frac{1}{\sigma_\alpha}s_t
\end{aligned}$$

Por lo que respecto al sector exterior de nuestra economía, tendríamos que:

$$\begin{aligned}
    &\text{Balanza comercial}:&& nx_t = y_t - c_t - \alpha s_t\\[0.5em]
    &&&\text{Combinando con (1)}\Longrightarrow \boxed{nx_t = \alpha\left(\frac{\omega}{\sigma} - 1\right)\, s_t}
\end{aligned}$$

De nuevo, en el caso especial $\sigma = \eta = \gamma = 1$, se tiene $nx_t = 0$ para todo $t$, aunque esta propiedad también se cumple para cualquier configuración de parámetros que satisfaga $\alpha(\gamma - 1) + (1 - \alpha)(\sigma\eta - 1) = 0$. Más en general, el signo de la relación entre los términos de intercambio y las exportaciones netas es ambiguo, dependiendo del tamaño relativo de $\alpha, \gamma$ y $\eta$.\footnote{%
    El hecho de que en nuestra economía se permitan variaciones en la balanza comercial constituye una diferencia clave respecto a muchos modelos de la literatura (véase, por ejemplo, Corsetti y Pesenti, 2001), que típicamente requieren utilidad logarítmica $(\sigma = 1)$ y elasticidad de sustitución unitaria $(\eta = 1)$ para que el equilibrio comercial se mantenga continuamente. Obsérvese también que nuestro marco exige condiciones más estrictas para el equilibrio comercial, ya que requiere además $\gamma = 1$ (o cualquier combinación de $\sigma$ y $\gamma$ tal que $\sigma \gamma = 1$).
}

Por último (la curva de PHILLIPS), en la pequeña economía abierta la dinámica de la inflación doméstica en términos del coste marginal real viene descrita por una ecuación análoga a la de una economía cerrada: $\pi_{H,t} = \beta E_t\{\pi_{H,t+1}\} + \kappa\, \widehat{mc}_t$. Donde los costes marginales son en términos reales y $\kappa$ es función del nivel de rigideces que haya en la economía ($\theta$) y $\beta$.

No obstante, la determinación del coste marginal real como función de la producción doméstica en una pequeña economía abierta difiere de la economía cerrada, debido a la existencia de una cuña entre producción y consumo, y entre precios domésticos y precios al consumo:

$$\begin{aligned}
    &\text{Curva de PHILLIPS}:\qquad \pi_{H,t} = \beta E_t\{\pi_{H,t+1}\} + \kappa\, \widehat{mc}_t\\[2em]
    &\text{Y costes marginales reales en  función de ($y_t,a_t, y_t^*$)}:\\
    &\Longrightarrow
    \left\{\begin{aligned}
        &(1)\quad mc_t = -\mu + (w_t - p_{H,t}) - a_t\\[0.5em]
        &(2)\quad\text{Utilizando}:
        \left\{\begin{aligned}
            &y_t=a_t+n_t\\
            &c_t=c_t^*+\left(\frac{1-\alpha}{\sigma}\right)s_t\\
            &\hspace{2em}\text{y,}\\
            &y_t=y_t^*+\frac{1}{\sigma_\alpha}_t
        \end{aligned}\right.
    \end{aligned}\right\}\Longrightarrow mc_t=-\nu+(\sigma_\alpha+\varphi)y_t+(\sigma-\sigma_\alpha)y_t^*-(1+\varphi)a_t
\end{aligned}$$

De este modo, el coste marginal aumenta con los términos de intercambio y con la producción mundial. Ambas variables influyen en el salario real a través del efecto riqueza sobre la oferta de trabajo, derivado de su impacto sobre el consumo doméstico. Además, los cambios en los términos de intercambio afectan directamente al salario real en términos del producto (product wage), para un salario real dado. La influencia de la tecnología (a través de su efecto directo sobre la productividad del trabajo) y de la producción doméstica (a través de su efecto sobre el empleo y, por tanto, sobre el salario real, dado el nivel de producción) es análoga a la observada en la economía cerrada.

Obsérvese que, en los casos especiales $\alpha = 0$ y/o $\sigma = \eta = \gamma = 1$, que implican $\sigma = \sigma\alpha$, el coste marginal real doméstico queda completamente aislado de las variaciones en la producción extranjera.

Por ultima, la representación canónica que utilizaremos en nuestra exposición. Se trata de la dinámica de equilibrio linealizada de una pequeña economía abierta, que puede representarse en términos de la brecha de producción ($x_t$) y la inflación doméstica ($\pi_{H,t}$) de forma análoga a su equivalente en una economía cerrada.\footnote{%
    El nivel natural de producción doméstica se obtiene imponiendo $mc_t = -\mu$ para todo $t$ y resolviendo para $y_t$: $\bar y_t = \Omega + \Gamma a_t + \alpha\Psi y_t^*$.
}

El coste marginal real doméstico y la brecha de producción están relacionados según $\widehat{mc}_t = (\sigma_\alpha + \varphi)\, x_t$, lo cual permite derivar una versión de la curva de Phillips neokeynesiana (NKPC) para la economía abierta en términos de la brecha de producción.

$$\begin{aligned}
    &\text{Curva de PHILLIPS}:&&\boxed{\pi_{H,t} = \beta E_t\{\pi_{H,t+1}\} + \kappa_x x_t}\\
    &\text{Curva IS}:&&\boxed{x_t = E_t\{x_{t+1}\} - \frac{1}{\sigma\psi}\big(r_t - E_t\{\pi_{H,t+1}\} - r_t^n\big)}
\end{aligned}$$

En general, la forma de la ecuación de PHILLIPS en la economía abierta es análoga a la de la economía cerrada en lo que respecta a la inflación doméstica. \textbf{El grado de apertura $\alpha$ afecta a la dinámica de la inflación únicamente a través de su impacto sobre la pendiente de la curva de PHILLIPS, es decir, sobre la sensibilidad de la inflación ante variaciones en la brecha de producción}. En una economía abierta, cambios en el output doméstico afectan al coste marginal tanto a través del empleo (capturado por $\varphi$) como de los términos de intercambio (capturados por $\sigma_\alpha$). En particular, si $\sigma\eta > 1$, una mayor apertura reduce el ajuste necesario de los términos de intercambio para absorber cambios en el output doméstico (relativo al mundial), amortiguando así su impacto sobre el coste marginal y la inflación.

Respecto a la Curva IS, esta se expresa para la economía abierta en términos de la brecha de producción. Así, el equilibrio se caracteriza por una ecuación tipo IS forward looking, similar a la de la economía cerrada. Sin embargo, hay dos diferencias importantes:
\begin{lnum}
    \item En primer lugar, el grado de apertura afecta a la sensibilidad de la brecha de producción a los tipos de interés. En particular, si $\omega > 1$ (que se obtiene con valores \textit{altos} de $\eta$ y $\gamma$), una mayor apertura aumenta dicha sensibilidad, debido al papel amplificador de los términos de intercambio sobre la demanda;
    \item En segundo lugar, la apertura introduce dependencia del tipo de interés natural respecto al crecimiento esperado del output mundial, además de la productividad doméstica.
\end{lnum}






\end{document}